% ---------------------------------------------------------------------------
% Notatki do kursu "Własny procesor RISC-V w SystemVerilogu"
% Budowanie: make   (pdflatex przez latexmk)
% ---------------------------------------------------------------------------
\documentclass[11pt,a4paper,oneside,parskip=half,headings=standardclasses,
  chapterprefix=true,numbers=noenddot]{scrbook}
\usepackage{kurs}

\hypersetup{pdftitle={Własny procesor RISC-V w SystemVerilogu --- notatki},
  pdfauthor={Kurs CPU Verilog}}

\begin{document}
\frontmatter
\begin{titlepage}
  \centering
  \vspace*{3cm}
  {\sffamily\color{szary}\large NOTATKI DO KURSU\par}
  \vspace{0.6cm}
  {\sffamily\bfseries\color{akcent}\fontsize{30}{36}\selectfont
    Własny procesor RISC-V\\[4pt] w SystemVerilogu\par}
  \vspace{0.8cm}
  {\large Od bramek logicznych do 5-etapowego potoku,\\
    na którym uruchamiasz programy w C\par}
  \vspace{2cm}
  \begin{tikzpicture}[scale=0.9, every node/.style={font=\sffamily\scriptsize}]
    \foreach \n/\l [count=\i] in {IF/pobierz, ID/dekoduj, EX/wykonaj, MEM/pamięć, WB/zapisz} {
      \node[blok, minimum width=18mm, minimum height=11mm] (s\i) at ({2.6*\i},0) {\textbf{\n}\\\l};
    }
    \foreach \i [evaluate=\i as \j using int(\i+1)] in {1,...,4} {
      \draw[->, thick, akcent] (s\i) -- (s\j);
    }
  \end{tikzpicture}
  \vfill
  {\sffamily\small\color{szary} Kod kursu: \texttt{\textasciitilde/Pulpit/kurs-cpu-verilog}\\
  Wersja z \today\par}
\end{titlepage}

\tableofcontents

\mainmatter
\addchap{Wstęp}

\section*{O czym jest ten kurs}

Celem kursu jest zbudowanie od zera procesora RISC-V (RV32I) w SystemVerilogu,
najpierw w wersji jednocyklowej, a potem jako klasyczny 5-etapowy potok,
i uruchomienie na nim programów skompilowanych z C. Po drodze powstaje
wszystko, czego potrzebuje projektant procesora: emulator referencyjny,
infrastruktura testowa, co-symulacja i fuzzing.

Notatki zakładają, że znasz C i programowanie niskopoziomowe (wskaźniki,
pamięć, ABI, asembler w zarysie). Nie zakładają znajomości języków opisu
sprzętu ani elektroniki. Tam, gdzie intuicja programisty C prowadzi na
manowce, znajdziesz czerwone ramki \emph{Pułapka}.

\section*{Jak czytać te notatki}

Każdy rozdział odpowiada jednemu modułowi kursu (katalogowi \plik{NN-nazwa/}).
Rozdział zaczyna się od metryczki (cel, czas, katalog), potem jest teoria z
przykładami, zadania (zielone ramki) i pytania kontrolne. Rozwiązania zadań
piszesz w plikach \plik{rtl/*.sv} danego modułu, a gotowe testy w
\plik{tb/} mówią, czy działa.

\begin{center}
\small
\begin{tabularx}{\linewidth}{@{}c l X r@{}}
\toprule
\textbf{Rozdz.} & \textbf{Moduł} & \textbf{Co budujesz} & \textbf{Wieczory} \\
\midrule
0 & \plik{00-start} & narzędzia, licznik, pierwszy testbench & 1 \\
1 & \plik{01-kombinacyjna} & mux, enkoder, sumator, barrel shifter; szerokości, zatrzaski & 2 \\
2 & \plik{02-sekwencyjna} & rejestry, timing, pamięci, FSM, FIFO, UART & 2--3 \\
3 & \plik{03-alu-regfile} & ALU, bank rejestrów, komparator skoków & 1--2 \\
4 & \plik{04-isa-emulator} & RV32I na poziomie bitów, generator immediate'ów, emulator w C & 2--3 \\
5 & \plik{05-single-cycle} & dekoder, LSU, procesor jednocyklowy & 3--4 \\
6 & \plik{06-weryfikacja} & co-symulacja, fuzzing, Verilator, mutacje & 1--2 \\
7 & \plik{07-c-bare-metal} & crt0, skrypt linkera, \code{\_\_mulsi3}: C na Twoim CPU & 1--2 \\
8 & \plik{08-pipeline} & potok 5-etapowy: forwarding, stalle, flush & 4--6 \\
9 & \plik{09-dalej} & rozszerzenie M, przerwania, FPGA, cache, formalna & --- \\
\bottomrule
\end{tabularx}
\end{center}

Moduły trzeba robić po kolei. Moduł 5 używa Twoich plików z modułów 3 i 4,
moduł 6 używa emulatora z modułu 4, a moduł 8 używa wszystkiego.

\section*{Narzędzia}

\begin{tabularx}{\linewidth}{@{}l X@{}}
\toprule
\textbf{Narzędzie} & \textbf{Do czego} \\
\midrule
Icarus Verilog (\code{iverilog}, \code{vvp}) & symulacja zdarzeniowa, 4 wartości logiczne, główny symulator kursu \\
GTKWave & oglądanie przebiegów (pliki \code{.vcd}) \\
Verilator & lint (\emph{-Wall dla sprzętu}) i szybka symulacja \\
Yosys & synteza: ile bramek, jak długa ścieżka krytyczna \\
\code{rvtool.py} & asembler i disasembler RV32I(M) dołączony do kursu \\
clang + \code{ld.lld} albo GCC dla RISC-V & tylko w rozdziale 7 (C na procesorze) \\
\bottomrule
\end{tabularx}

\medskip
Do rozdziału 7 potrzebny jest linker: \code{sudo apt install lld} (albo
pełny toolchain \code{gcc-riscv64-unknown-elf}). Pozostałe rozdziały nie
wymagają żadnego toolchaina RISC-V.

\section*{Rytm pracy}

W każdym module te same polecenia (uruchamiane z katalogu modułu):

\begin{sh}
make test               # wszystkie testy modułu (na starcie: same FAIL)
make unit T=mux4        # jeden test jednostkowy
make wave T=mux4        # przebiegi w GTKWave
make lint               # Verilator --lint-only
make synth T=adder      # Yosys: bramki i najdłuższa ścieżka
\end{sh}

W modułach z procesorem (5--8) dochodzą:

\begin{sh}
make progs              # 12 programów testowych na Twoim rdzeniu
make prog P=08_fib      # jeden program: wyjście + ślad wykonania
make cosim              # porównanie śladu z emulatorem
make fuzz N=100         # losowe programy + porównanie z emulatorem
make progs SIM=verilator
\end{sh}

\begin{uwaga}[Rada]
Zrób z katalogu kursu repozytorium git i commituj po każdym zadaniu. Przy
potoku (rozdział 8) możliwość powrotu do ostatniej działającej wersji jest
bezcenna.
\end{uwaga}

\section*{Konwencje}

\begin{itemize}
  \item SystemVerilog w podzbiorze rozumianym przez Icarus 12, Verilatora i
    Yosysa: \code{logic}, \code{always\_comb}, \code{always\_ff},
    \code{typedef enum}, pakiety.
  \item \code{always\_comb} z przypisaniem \code{=} dla logiki kombinacyjnej,
    \code{always\_ff} z \code{<=} dla rejestrów. Bez wyjątków.
  \item Jeden zegar \code{clk}, reset \emph{synchroniczny}, aktywny w stanie
    wysokim (\code{rst}).
  \item Stałe architektury zawsze z pakietu \code{rv32i\_pkg}
    (\plik{common/rtl/rv32i\_pkg.sv}), nigdy „magiczne liczby”.
\end{itemize}

\setcounter{chapter}{-1}
\chapter{Narzędzia i pierwszy układ}
\label{ch:start}

\metryczka{zrozumieć, czym opis sprzętu różni się od programu, i opanować cykl
\emph{napisz → zasymuluj → obejrzyj przebiegi → zlintuj → zsyntezuj}.}
{1 wieczór}{00-start/}

\section{HDL to nie jest język programowania}

W C piszesz \emph{sekwencję kroków} dla procesora, który już istnieje. W
SystemVerilogu opisujesz \emph{układ}: zbiór bramek i przerzutników
połączonych przewodami. Wszystkie elementy działają \emph{jednocześnie} i
\emph{bez przerwy}.

\begin{sv}
assign y = a & b;   // to NIE jest "oblicz a&b i zapisz do y":
                    // to bramka AND na stałe połączona z przewodem y
\end{sv}

Ten sam kod ma dwa życia:

\begin{center}
\begin{tabularx}{0.92\linewidth}{@{}l X X@{}}
\toprule
& \textbf{Symulacja} & \textbf{Synteza} \\
\midrule
Narzędzie & Icarus Verilog, Verilator & Yosys (dalej: nextpnr → FPGA) \\
Co robi & wykonuje opis jak program zdarzeniowy & zamienia opis na sieć bramek i przerzutników \\
Co akceptuje & prawie wszystko (\code{\#5}, \code{\$display}, pliki) & tylko \emph{syntezowalny podzbiór} \\
\bottomrule
\end{tabularx}
\end{center}

Testbench (\plik{tb/*.sv}) żyje tylko w symulacji i może używać
wszystkiego. Kod w \plik{rtl/} musi dać się zsyntezować, co w każdej
chwili sprawdzisz poleceniem \code{make synth T=moduł}.

\begin{pulapka}[Pułapka dla programisty C]
Kod, który się symuluje, wcale nie musi się syntezować do tego, co myślisz.
Pętla \code{for} to $N$ \emph{kopii} sprzętu, a nie $N$ iteracji w czasie.
\code{if} bez \code{else} w logice kombinacyjnej tworzy zatrzask. Operatory
\code{/} i \code{\%} budują ogromny dzielnik. Zawsze pytaj: \emph{jak
wyglądałby ten układ na schemacie?}
\end{pulapka}

\section{Dwa rodzaje logiki}

\paragraph{Logika kombinacyjna:} wyjście zależy \emph{tylko} od bieżących
wejść, bez pamięci. Bramki, multipleksery, sumatory, dekodery.
\begin{sv}
always_comb y = sel ? a : b;      // multiplekser
assign sum = a + b;               // sumator
\end{sv}

\paragraph{Logika sekwencyjna:} przerzutniki (\emph{flip-flopy})
zapamiętują wartość na zboczu zegara. To jest \emph{stan} układu.
\begin{sv}
always_ff @(posedge clk) q <= d;  // rejestr: na zboczu narastającym q := d
\end{sv}

Prawie każdy układ synchroniczny ma tę samą budowę (rys.~\ref{fig:sync}):
logika kombinacyjna liczy \emph{następny stan} z bieżącego stanu i
wejść, a rejestry na zboczu zegara \emph{jednocześnie} go zapamiętują.

\begin{figure}[h]
\centering
\begin{tikzpicture}[node distance=12mm]
  \node[blok, minimum width=32mm, minimum height=14mm] (comb) {logika\\kombinacyjna};
  \node[blok, fill=zadaniejasny, right=18mm of comb, minimum height=14mm] (reg) {rejestry\\(przerzutniki)};
  \draw[->, thick] ($(comb.west)+(-16mm,3mm)$) node[left, etykieta]{wejścia} -- ($(comb.west)+(0,3mm)$);
  \draw[->, thick] (comb) -- node[above, sygnal]{d (następny)} (reg);
  \draw[->, thick] (reg.east) -- ++(14mm,0) node[right, etykieta]{wyjścia};
  \draw[->, thick] ($(reg.east)+(6mm,0)$) |- ($(comb.south)+(0,-9mm)$) -- node[left, sygnal, pos=0.3]{q (stan)} ($(comb.south)+(0,0)$);
  \draw[->, thick] ($(reg.south)+(0,-12mm)$) node[below, sygnal]{clk} -- (reg.south);
\end{tikzpicture}
\caption{Ogólny model układu synchronicznego.}
\label{fig:sync}
\end{figure}

\section{Anatomia modułu}

Plik \plik{rtl/counter.sv} (gotowy przykład):

\begin{sv}
module counter #(
  parameter int WIDTH = 8           // parametr: stała ustalana przy tworzeniu instancji
) (
  input  logic             clk,
  input  logic             rst,
  input  logic             en,
  output logic [WIDTH-1:0] q        // wektor WIDTH bitów, bit 0 najmłodszy
);
  always_ff @(posedge clk) begin
    if (rst)     q <= '0;
    else if (en) q <= q + 1'b1;
  end
endmodule
\end{sv}

\begin{itemize}
  \item \code{logic} to typ SystemVerilog dla sygnałów, zastępujący stare
    \code{wire} i \code{reg}. W kursie używamy wyłącznie \code{logic}.
  \item \code{[WIDTH-1:0]} to zakres bitów: \code{q[0]} najmłodszy,
    \code{q[WIDTH-1]} najstarszy.
  \item Literały mają postać \code{<szerokość>'<baza><wartość>}, np.
    \code{8'hFF}, \code{4'b1010}, \code{32'd100}, \code{1'b1}. Zapisy
    \code{'0} i \code{'1} oznaczają „same zera” i „same jedynki” w
    dowolnej szerokości.
  \item Reset jest \emph{synchroniczny} (sprawdzany na zboczu) i aktywny w
    stanie 1. Tej konwencji trzymamy się w całym kursie.
\end{itemize}

Instancja, czyli „wstawienie układu do innego układu”:
\begin{sv}
counter #(.WIDTH(16)) u_cnt (.clk(clk), .rst(rst), .en(1'b1), .q(value));
counter #(.WIDTH(16)) u_cnt (.clk, .rst, .en(1'b1), .q(value)); // .clk = .clk(clk)
\end{sv}

\section{Anatomia testbencha}

Testbench to moduł bez portów, który tworzy sygnały, wstawia testowany układ
(DUT, \emph{device under test}), generuje zegar i w bloku \code{initial}
podaje pobudzenia oraz sprawdza wyniki:

\begin{sv}
module tb_counter;
  `include "tb_common.svh"         // makra CHECK_EQ, CHECK, tb_start/tb_finish
  logic clk = 0, rst, en;
  logic [7:0] q;
  counter #(.WIDTH(8)) dut (.clk, .rst, .en, .q);
  always #5 clk = ~clk;            // okres 10 jednostek czasu

  initial begin
    tb_start();
    rst = 1; en = 0;
    @(negedge clk);                // czekaj na zbocze OPADAJĄCE
    @(negedge clk);
    `CHECK_EQ(q, 8'd0, "po resecie")
    rst = 0; en = 1;
    repeat (5) @(negedge clk);
    `CHECK_EQ(q, 8'd5, "po 5 taktach liczenia")
    tb_finish();
  end
endmodule
\end{sv}

\subsection*{Dlaczego wejścia zmieniamy na zboczu opadającym?}

Przerzutniki reagują na zbocze narastające. Gdyby testbench zmieniał wejście
\emph{dokładnie} w chwili zbocza narastającego, wynik zależałby od tego, w
jakiej kolejności symulator przetworzy oba zdarzenia. To wyścig
(\emph{race}). Zmiana w połowie okresu usuwa problem (rys.~\ref{fig:cnt}).

\begin{figure}[h]
\centering
\begin{tikzpicture}[x=6mm, y=6mm]
  \node[sygnal, anchor=east] at (-0.3,0.25) {clk};
  \foreach \i in {0,...,8} { \draw[thick] ({2*\i},0) -- ++(0,0.5) -- ++(1,0) -- ++(0,-0.5) -- ++(1,0); }
  \node[sygnal, anchor=east] at (-0.3,-1.25) {rst};
  \draw[thick] (0,-1) -- (3,-1) -- (3,-1.5) -- (18,-1.5);
  \node[sygnal, anchor=east] at (-0.3,-2.75) {en};
  \draw[thick] (0,-3) -- (3,-3) -- (3,-2.5) -- (13,-2.5) -- (13,-3) -- (18,-3);
  \node[sygnal, anchor=east] at (-0.3,-4.25) {q};
  \draw[thick] (0,-4) -- (2,-4) (0,-4.5) -- (2,-4.5);
  \node[sygnal] at (1,-4.25) {X};
  \foreach \a/\b/\v in {2/4/0, 4/6/1, 6/8/2, 8/10/3, 10/12/4, 12/18/5} {
    \draw[thick] (\a,-4) -- (\b,-4) (\a,-4.5) -- (\b,-4.5);
    \draw[thick] (\a,-4) -- (\a,-4.5);
    \node[sygnal] at ({(\a+\b)/2},-4.25) {\v};
  }
  \foreach \x in {3,5,7,9,11,13,15} { \draw[densely dotted, szary] (\x,0.8) -- (\x,-4.8); }
  \node[etykieta, align=center] at (9,-5.6) {kropki: zbocza opadające --- tu testbench zmienia wejścia i sprawdza wyniki};
\end{tikzpicture}
\caption{Przebiegi testu licznika: \code{q} zmienia się tuż po zboczu narastającym,
na którym \code{en = 1} (a przed resetem ma wartość nieznaną X).}
\label{fig:cnt}
\end{figure}

\section{Logika czterowartościowa: 0, 1, X, Z}

Symulator zna cztery wartości bitu: \code{0}, \code{1}, \code{X}
(\emph{nieznana}: niezainicjowany przerzutnik, konflikt dwóch sterowników
albo wynik operacji na X) oraz \code{Z} (nikt nie steruje przewodem; w kursie
nie występuje).

X jest Twoim sprzymierzeńcem. Czerwone \code{X} w GTKWave oznacza coś
niezainicjowanego albo nieprzypisanego. Prawdziwy sprzęt miałby tam losowe 0
lub 1 i błąd ujawniałby się „czasem”, a symulacja pokazuje go od razu.

\begin{pulapka}[Porównania z X]
\code{a == b} z X daje X, co w \code{if} działa jak \emph{fałsz}! Operatory
\code{===} i \code{!==} porównują dosłownie, łącznie z X. W testbenchach
używaj \code{===}/\code{!==} (robi to makro \code{CHECK\_EQ}), w RTL zwykłego
\code{==}.
\end{pulapka}

\section{Jak symulator wykonuje kod}

Symulator zdarzeniowy utrzymuje kolejkę zdarzeń uporządkowaną w czasie. W
jednej chwili czasu najpierw wykonuje wszystkie procesy, które się
„obudziły” (zmiana sygnału z listy wrażliwości, zbocze zegara), a przypisania
nieblokujące \code{<=} odkłada na koniec tej chwili (region NBA). Dlatego
wszystkie \code{always\_ff} widzą \emph{stare} wartości rejestrów, tak jak
prawdziwe przerzutniki taktowane jednym zegarem. Wrócimy do tego w
rozdziale~\ref{ch:sekw}.

\section{GTKWave w minutę}

\begin{enumerate}
  \item \code{make wave T=updown} otwiera GTKWave z przebiegami testu.
  \item Po lewej, w drzewku \emph{SST}, rozwiń \code{tb\_updown}, potem \code{dut}.
  \item Zaznacz sygnały i kliknij \emph{Append} (albo przeciągnij je na panel).
  \item \code{Ctrl+Shift+F}: dopasuj widok do całego czasu symulacji.
  \item Prawy klik na wektorze, \emph{Data Format}: \emph{Decimal} albo \emph{Hex}.
  \item \emph{File → Write Save File}: zapamiętuje układ sygnałów.
\end{enumerate}

\section{Zadania}

\begin{zadanie}{0.1 --- Uruchom gotowy licznik}
\code{make unit T=counter}, potem \code{make wave T=counter}. Dodaj
\code{clk}, \code{rst}, \code{en}, \code{q}. Znajdź moment, w którym \code{en}
spada do 0, i sprawdź, czy \code{q} się zatrzymuje. W którym dokładnie
momencie zmienia się \code{q} i względem którego zbocza?
\end{zadanie}

\begin{zadanie}{0.2 --- Licznik w górę/w dół (\code{rtl/updown.sv})}
Priorytety: \code{rst} > \code{load} > \code{en}. Przy \code{en = 1} licznik
liczy w górę albo w dół (\code{up}) z zawijaniem. Wyjście kombinacyjne
\code{tc} (\emph{terminal count}) = 1 dokładnie wtedy, gdy \code{en = 1} i
następny krok spowoduje zawinięcie. Test porównuje układ z modelem przez 500
losowych taktów i wypisuje pierwsze rozbieżności z czasem symulacji. Użyj go,
żeby znaleźć to miejsce w GTKWave.
\end{zadanie}

\begin{zadanie}{0.3 --- Pierwsza synteza}
\code{make synth T=counter}. Yosys wypisze m.in. 8 komórek
\code{\$\_SDFFE\_PP0P\_}: przerzutnik z synchronicznym resetem do 0 i
zezwoleniem, czyli dokładnie to, co opisałeś. Zmień domyślne
\code{WIDTH} na 32 i porównaj liczbę bramek oraz najdłuższą ścieżkę. Dlaczego
ścieżka rośnie? (Podpowiedź: przeniesienie w dodawaniu.)
\end{zadanie}

\begin{zadanie}{0.4 --- Lint}
\code{make lint} powinno zgłosić \code{OK}. Dla eksperymentu przypisz
\code{q <= d[1:0];} przy \code{WIDTH = 4} i zobacz ostrzeżenie o
szerokości. Verilator to Twój \code{-Wall -Werror}.
\end{zadanie}

\begin{gotowe}
\code{make test} pokazuje PASS dla \code{counter} i \code{updown}, a \code{make lint} nie zgłasza ostrzeżeń.
\end{gotowe}

\begin{kontrolne}
  \item Czym różni się \code{assign y = a + b;} od \code{y = a + b;} w funkcji C?
  \item Dlaczego testbench zmienia wejścia na zboczu opadającym zegara?
  \item Co oznacza X w przebiegach i czemu \code{if (x == 1)} przy \code{x = X} wybiera gałąź \code{else}?
  \item Które konstrukcje z \plik{tb\_counter.sv} nie dałyby się zsyntezować i dlaczego to nie problem?
\end{kontrolne}

\chapter{Logika kombinacyjna}
\label{ch:komb}

\metryczka{pisać logikę kombinacyjną świadomie: wiedzieć, jaki sprzęt powstaje z
każdej linii, i nie dać się złapać na szerokości wyrażeń ani zatrzaski.}
{2 wieczory}{01-kombinacyjna/}

\section{\texttt{assign} i \texttt{always\_comb}}

Obie konstrukcje opisują logikę kombinacyjną:

\begin{sv}
assign y = sel ? a : b;

always_comb begin
  if (sel) y = a;
  else     y = b;
end
\end{sv}

W \code{always\_comb} używamy przypisania \emph{blokującego} \code{=}. Wewnątrz
bloku działa jak w C, bo kolejne linie widzą wyniki poprzednich, ale cały blok to
\emph{jedna} sieć bramek liczona „natychmiast”. Zmienna pośrednia jest tylko nazwą
dla przewodu:

\begin{sv}
always_comb begin
  t = a ^ b;        // t: nazwa dla przewodu z wyjścia bramki XOR
  y = t & c;
end
\end{sv}

\code{always\_comb} (w odróżnieniu od starego \code{always @(*)}) informuje też
narzędzia o Twoim zamiarze: jeśli wyjdzie z niego coś, co nie jest
kombinacyjne (zatrzask), Yosys odmówi syntezy.

\section{Wektory: wycinanie, sklejanie, powielanie}

\begin{sv}
logic [31:0] x;
x[7:0]                    // najmłodszy bajt
x[31]                     // bit znaku
x[8*i +: 8]               // bajt i: 8 bitów od pozycji 8*i w górę (i może być zmienną)
{a, b}                    // sklejenie: a na starszych bitach
{4{b}}                    // b powtórzone 4 razy
{{20{x[31]}}, x[31:20]}   // rozszerzenie znakiem 12 -> 32 bity (rozdział 4!)
{cout, sum} = a + b;      // sklejenie po lewej: wynik z przeniesieniem
\end{sv}

\section{Liczby ze znakiem w sprzęcie}

Sprzęt nie „wie”, czy wektor bitów jest liczbą ze znakiem. Kod
uzupełnień do dwóch (U2) jest wybrany właśnie dlatego, że \emph{ten sam
sumator} poprawnie dodaje liczby ze znakiem i bez znaku. Znak ma znaczenie
tylko w kilku operacjach:

\begin{center}
\begin{tabular}{@{}lll@{}}
\toprule
Operacja & bez znaku & ze znakiem \\
\midrule
porównanie \code{<} & \code{a < b} & \code{\$signed(a) < \$signed(b)} \\
przesunięcie w prawo & \code{a >> n} (wsuwa 0) & \code{\$signed(a) >>> n} (powiela bit znaku) \\
rozszerzanie do szerszego typu & zerami & kopiami bitu znaku \\
mnożenie (górna połowa), dzielenie & osobne układy & osobne układy \\
\bottomrule
\end{tabular}
\end{center}

Negacja w U2: $-x = \overline{x} + 1$. Stąd odejmowanie $a - b = a +
\overline{b} + 1$, czyli \emph{sumator z negacją jednego wejścia i
przeniesieniem wejściowym 1}. Z tego triku skorzysta ALU w rozdziale 3.

\section{Szerokości wyrażeń: najczęstsze źródło błędów}

W C typ wyrażenia zależy od typów argumentów. W SystemVerilogu \textbf{szerokość
obliczeń zależy też od lewej strony przypisania}:

\begin{sv}
logic [7:0] a, b, s8;
logic [8:0] s9;
s8 = a + b;   // liczone na 8 bitach: przeniesienie ginie
s9 = a + b;   // liczone na 9 bitach: przeniesienie trafia do s9[8]
\end{sv}

\begin{itemize}
  \item Wyrażenie arytmetyczne liczy się w szerokości $\max$(operandy, lewa
    strona). Węższe operandy są rozszerzane.
  \item Liczba bez rozmiaru (\code{5}) ma 32 bity i jest \emph{signed}.
    \code{1'b1} ma 1 bit.
  \item \textbf{Jeśli którykolwiek operand jest unsigned, całe wyrażenie jest
    unsigned.} \code{logic} jest domyślnie unsigned.
\end{itemize}

\begin{pulapka}[\texttt{>>>} bez \texttt{\$signed}]
\code{a >>> n} na wektorze unsigned to zwykłe \code{>>}. Przesunięcie
arytmetyczne wymaga \code{\$signed(a) >>> n}. Podobnie \code{a < b} na
\code{logic} porównuje bez znaku. Ten błąd wychodzi dopiero na liczbach ujemnych,
dlatego testy kursu celowo używają \code{0x80000000} i \code{0xFFFFFFFF}.
\end{pulapka}

Verilator (\code{make lint}) ostrzega o niejawnym obcinaniu i rozszerzaniu
(\code{WIDTHTRUNC}, \code{WIDTHEXPAND}). Traktuj te ostrzeżenia poważnie.

\section{Jaki sprzęt powstaje}

\begin{center}
\begin{tabularx}{\linewidth}{@{}l X l@{}}
\toprule
\textbf{Konstrukcja} & \textbf{Sprzęt} & \textbf{Koszt} \\
\midrule
\code{\& | \^{} \~{}} & bramki, po jednej na bit & 1 poziom \\
\code{c ? a : b}, \code{if/else} & multiplekser 2:1 & tani \\
\code{case} na $N$ wartości & multiplekser $N$:1 (drzewo, $\log_2 N$ poziomów) & umiarkowany \\
\code{a == b} & XOR-y i drzewo OR & tani \\
\code{a + b}, \code{a - b}, \code{a < b} & sumator (łańcuch przeniesień) & średni \\
\code{a << b} ($b$ zmienne) & $\log_2 N$ warstw multiplekserów & spory \\
\code{a * b} & macierz sumatorów & \textbf{drogi} \\
\code{a / b}, \code{a \% b} & sieć odejmowań & \textbf{bardzo drogi} \\
\bottomrule
\end{tabularx}
\end{center}

\paragraph{Ścieżka krytyczna.} Sygnał musi przejść przez najdłuższy łańcuch
bramek w ciągu jednego taktu zegara. Im dłuższy łańcuch, tym niższa
maksymalna częstotliwość. \code{make synth} podaje to jako \emph{longest
topological path}, czyli liczbę bramek na najdłuższej drodze. To przybliżenie, bo
prawdziwe opóźnienia zależą od technologii, ale dobrze pokazuje trendy.

\section{Przykład: sumator ripple-carry}

Sumator pełny (\emph{full adder}) dodaje trzy bity:
\[ s = a \oplus b \oplus c_{in}, \qquad c_{out} = ab + c_{in}(a \oplus b). \]
Łącząc $N$ takich sumatorów łańcuchem przeniesień, dostajemy sumator
$N$-bitowy (rys.~\ref{fig:rca}). Przeniesienie z bitu 0 może wpłynąć na bit
$N-1$, więc ścieżka krytyczna jest proporcjonalna do $N$: $O(n)$.

\begin{figure}[h]
\centering
\begin{tikzpicture}[x=24mm]
  \foreach \i/\n in {0/0, 1/1, 2/2, 4/N\!-\!1} {
    \node[blok, minimum width=14mm, minimum height=10mm] (fa\i) at (-\i,0) {FA};
    \draw[<-, thick] ($(fa\i.north)+(-3mm,0)$) -- ++(0,6mm) node[above left=-1mm, sygnal]{a$_{\n}$};
    \draw[<-, thick] ($(fa\i.north)+(3mm,0)$) -- ++(0,6mm) node[above right=-1mm, sygnal]{b$_{\n}$};
    \draw[->, thick] (fa\i.south) -- ++(0,-6mm) node[below, sygnal]{s$_{\n}$};
  }
  \node at (-3,0) {$\cdots$};
  \draw[<-, thick] (fa0.east) -- ++(6mm,0) node[right, sygnal]{c$_{in}$};
  \draw[->, thick] (fa0.west) -- (fa1.east);
  \draw[->, thick] (fa1.west) -- (fa2.east);
  \draw[->, thick] (fa2.west) -- ($(-3,0)+(5mm,0)$);
  \draw[->, thick] ($(-3,0)+(-5mm,0)$) -- (fa4.east);
  \draw[->, thick] (fa4.west) -- ++(-6mm,0) node[left, sygnal]{c$_{out}$};
  \draw[decorate, decoration={brace, amplitude=5pt, mirror}, akcent] ($(fa4.south west)+(0,-11mm)$) -- node[below=6pt, etykieta]{ścieżka krytyczna: przez wszystkie przeniesienia} ($(fa0.south east)+(0,-11mm)$);
\end{tikzpicture}
\caption{Sumator ripple-carry z $N$ sumatorów pełnych.}
\label{fig:rca}
\end{figure}

Szybsze sumatory (carry-lookahead, Kogge--Stone) liczą przeniesienia
równolegle w $O(\log n)$ poziomach kosztem większej liczby bramek. Na FPGA
operator \code{+} trafia na dedykowane łańcuchy przeniesień, wielokrotnie
szybsze od sumatora złożonego ręcznie z LUT-ów. Wniosek praktyczny:
\textbf{pisz \code{+} i pozwól narzędziom wybrać implementację}. Ręczna
budowa ma sens tylko w celach dydaktycznych albo gdy wiesz, po co.

\section{\texttt{if} kontra \texttt{case} i priorytety}

\begin{sv}
always_comb begin            // łańcuch priorytetów: a wygrywa z b, b z c
  if (a)      y = 1;
  else if (b) y = 2;
  else if (c) y = 3;
  else        y = 0;
end
\end{sv}

Kolejne \code{else if} to kaskada multiplekserów (rys.~\ref{fig:prio}).
Warunek z \emph{pierwszego} \code{if} jest najbliżej wyjścia, więc ma najwyższy
priorytet. Gdy warunki się wykluczają (np. różne wartości jednego sygnału),
\code{case} lepiej oddaje intencję i daje syntezatorowi swobodę zbudowania
płaskiego drzewa.

\begin{figure}[h]
\centering
\begin{tikzpicture}[x=22mm]
  \foreach \i/\c/\v in {1/c/3, 2/b/2, 3/a/1} {
    \pic (m\i) at (\i,0) {muks};
    \draw[<-, thick] (m\i-sel) -- ++(0,-7mm) node[below, sygnal]{\c};
    \draw[<-, thick] (m\i-in1) -- ++(-7mm,0) node[left, sygnal]{\v};
  }
  \draw[<-, thick] (m1-in0) -- ++(-7mm,0) node[left, sygnal]{0};
  \draw[->, thick] (m1-out) -- ++(5mm,0) |- (m2-in0);
  \draw[->, thick] (m2-out) -- ++(5mm,0) |- (m3-in0);
  \draw[->, thick] (m3-out) -- ++(8mm,0) node[right, sygnal]{y};
\end{tikzpicture}
\caption{\code{if / else if / else} jako łańcuch multiplekserów 2:1. Wejście 1 jest wybierane, gdy warunek (nóżka u dołu) jest prawdziwy; warunek \code{a} jest najbliżej wyjścia, więc ma najwyższy priorytet.}
\label{fig:prio}
\end{figure}

\section{Zatrzaski: błąd numer jeden}

\begin{sv}
always_comb begin
  if (en) y = a;    // a gdy en == 0? y ma "pamiętać" starą wartość!
end
\end{sv}

Skoro \code{y} ma pamiętać wartość przy \code{en = 0}, syntezator musi wstawić
element pamiętający, czyli \emph{zatrzask} sterowany poziomem. W układach
synchronicznych to prawie zawsze błąd: psuje analizę czasową, a układ zachowuje
się inaczej na krzemie niż w symulacji.

\textbf{Lekarstwo:} każda zmienna przypisywana w \code{always\_comb} musi dostać
wartość na \emph{każdej} ścieżce. Najprościej przez przypisanie domyślne na
początku bloku:

\begin{sv}
always_comb begin
  y = '0;           // wartość domyślna
  if (en) y = a;
end
\end{sv}

To samo dotyczy \code{case} bez \code{default}, gdy nie wszystkie wartości są
pokryte.

\section{Pętle \texttt{for} i \texttt{generate}: kopiowanie sprzętu}

\begin{sv}
always_comb begin
  parity = 1'b0;
  for (int i = 0; i < 32; i++)
    parity = parity ^ x[i];       // 32 bramki XOR połączone łańcuchem
end
\end{sv}

Pętla jest \textbf{rozwijana w czasie syntezy}: nie ma tu licznika ani iteracji
w czasie, tylko 32 kopie logiki. Liczba iteracji musi być stała.
(Syntezator zwykle przebuduje ten łańcuch w drzewo o głębokości $\log_2 32$, bo XOR
jest łączny.)

Do powielania \emph{instancji} modułów służy \code{generate}:
\begin{sv}
genvar i;
generate
  for (i = 0; i < WIDTH; i++) begin : g_fa      // etykieta bloku jest wymagana
    full_adder fa (.a(a[i]), .b(b[i]), .cin(c[i]), .s(sum[i]), .cout(c[i+1]));
  end
endgenerate
\end{sv}

\section{Przykład: przesuwnik logarytmiczny}

Przesunięcie o zmienną liczbę bitów $0..31$ można zbudować z 5 warstw
multiplekserów: warstwa $k$ przesuwa o $2^k$ albo nie, zależnie od bitu
\code{shamt[k]} (rys.~\ref{fig:barrel}). Każda liczba $0..31$ jest sumą
niektórych potęg dwójki. Przesunięcie w lewo uzyskuje się, odwracając
kolejność bitów, przesuwając w prawo i odwracając z powrotem.

\begin{figure}[h]
\centering
\begin{tikzpicture}[x=19mm]
  \node[sygnal] (in) at (0,0) {a};
  \foreach \k/\s in {1/1, 2/2, 3/4, 4/8, 5/16} {
    \node[blok, minimum width=13mm, minimum height=9mm, font=\sffamily\scriptsize] (l\k) at (\k,0) {$\gg\!\s$\\lub nic};
    \draw[<-, thick] (l\k.south) -- ++(0,-5mm) node[below, sygnal]{shamt[\the\numexpr\k-1\relax]};
  }
  \draw[->, thick] (in) -- (l1);
  \foreach \k [evaluate=\k as \j using int(\k+1)] in {1,...,4} { \draw[->, thick] (l\k) -- (l\j); }
  \draw[->, thick] (l5) -- ++(13mm,0) node[right, sygnal]{y};
\end{tikzpicture}
\caption{Barrel shifter: 5 warstw po 32 multipleksery 2:1.}
\label{fig:barrel}
\end{figure}

\section{Znane ograniczenia Icarus Verilog 12}

Icarus nie obsługuje całego SystemVeriloga. W RTL rzadko to przeszkadza, ale
warto wiedzieć:
\begin{itemize}
  \item brak \code{break} i \code{continue} w pętlach,
  \item brak inicjalizacji całej tablicy w deklaracji (\code{logic [7:0] t [4] = '\{...\}}),
  \item \code{unique case} i \code{priority case} są akceptowane, ale ignorowane,
  \item komunikat \emph{sorry: constant selects in always\_* processes} to tylko
    ostrzeżenie (Makefile go ukrywa).
\end{itemize}

\section{Zadania}

\begin{zadanie}{1.1 --- Multiplekser (\code{rtl/mux4.sv})}
Dwie wersje: \code{always\_comb} z \code{case} oraz jeden \code{assign} z
operatorem \code{?:}. Test sprawdza instancje 8- i 16-bitową, czyli też to, czy
parametr działa.
\end{zadanie}

\begin{zadanie}{1.2 --- Enkoder priorytetowy (\code{rtl/prio\_enc.sv})}
\code{idx} = numer najstarszego ustawionego bitu, \code{valid} = czy jakikolwiek.
Użyj pętli \code{for}. Pętla od bitu 0 do 7, w której każde trafienie nadpisuje
\code{idx}, daje regułę „najstarszy wygrywa”. Dlaczego? Narysuj powstały
łańcuch multiplekserów.
\end{zadanie}

\begin{zadanie}{1.3 --- Sumator ripple-carry (\code{rtl/adder.sv})}
Najpierw \code{full\_adder} z samych bramek, potem \code{adder} przez
\code{generate}, bez operatora \code{+}. Test sprawdza \emph{wszystkie}
kombinacje dla 8 bitów ($2^{17}$ przypadków) i losowe dla 32. Potem
\code{make synth T=adder}: zanotuj liczbę komórek i ścieżkę, zastąp ciało
modułu wersją \code{assign \{cout, sum\} = a + b + cin;} i porównaj.
\end{zadanie}

\begin{zadanie}{1.4 --- Barrel shifter (\code{rtl/shifter.sv})}
SLL, SRL i SRA o \code{shamt} z 5 warstw multiplekserów, bez operatorów
przesunięcia. Ile multiplekserów 2:1 ma Twój układ (bez odwracania)? Porównaj z
\code{make synth T=shifter}.
\end{zadanie}

\begin{zadanie}{1.5 --- Polowanie na zatrzask (\code{rtl/hex7seg.sv})}
Dekoder 7-segmentowy ma dwa błędy prowadzące do zatrzasku. Zanim cokolwiek
poprawisz, uruchom \code{make lint}, \code{make synth T=hex7seg} i
\code{make unit T=hex7seg}. Zauważ, że lint znajduje tylko jeden z dwóch
błędów, a Yosys odmawia syntezy z komunikatem \emph{Latch inferred}. Potem
popraw kod i dokończ tablicę A--F.
\end{zadanie}

\begin{gotowe}
\code{make test} → 5 × PASS, \code{make lint} → wszystko OK.
\end{gotowe}

\begin{kontrolne}
  \item Ile bitów ma wynik \code{a + b} dla 8-bitowych \code{a}, \code{b} przy przypisaniu do 16-bitowej zmiennej? A do 8-bitowej?
  \item Co da \code{\$signed(8'hF0) >>> 2}, a co \code{8'hF0 >>> 2}?
  \item Dlaczego \code{if} bez \code{else} w \code{always\_comb} tworzy zatrzask, a w \code{always\_ff} nie?
  \item Czym różni się pętla \code{for} w RTL od pętli w C? Co jest stałe w czasie syntezy?
  \item Jak głęboka (w poziomach logiki) jest ścieżka krytyczna sumatora ripple-carry dla $N$ bitów, a jak przesuwnika logarytmicznego?
\end{kontrolne}

\chapter{Logika sekwencyjna i automaty stanów}
\label{ch:sekw}

\metryczka{rejestry, czas w układzie synchronicznym, pamięci i automaty
skończone. Procesor to w gruncie rzeczy wielki automat: jego stanem są PC,
rejestry i pamięć.}{2--3 wieczory}{02-sekwencyjna/}

\section{Przerzutnik i \texttt{always\_ff}}

\begin{sv}
always_ff @(posedge clk) q <= d;
\end{sv}

Na każdym zboczu narastającym \code{clk} przerzutnik zapamiętuje \code{d}.
Między zboczami \code{q} się nie zmienia, cokolwiek dzieje się z \code{d}.

\subsection*{Dlaczego przypisanie nieblokujące \texttt{<=}?}

\code{<=} oznacza: \emph{oblicz prawą stronę teraz, przypisz na końcu
kroku czasowego}. Wszystkie \code{<=} w całym układzie widzą więc \emph{stare}
wartości, dokładnie tak jak prawdziwe przerzutniki taktowane jednym zegarem:

\begin{sv}
always_ff @(posedge clk) begin   // rejestr przesuwny 3-bitowy
  a <= in;
  b <= a;                        // stare a
  c <= b;                        // stare b
end

always_ff @(posedge clk) begin   // zamiana wartości: działa!
  x <= y;
  y <= x;
end
\end{sv}

\begin{pulapka}[\texttt{=} w \texttt{always\_ff}]
Z \code{=} linia \code{b = a} widziałaby \emph{nowe} \code{a} i rejestr
przesuwny zapadłby się w jeden przerzutnik. Przy kilku blokach \code{always}
wynik zależałby od kolejności ich wykonania przez symulator. Zasada bez
wyjątków: \textbf{w \code{always\_ff} tylko \code{<=}}.
\end{pulapka}

\subsection*{Rejestr z zezwoleniem i resetem}

\begin{sv}
always_ff @(posedge clk) begin
  if (rst)     q <= '0;      // reset synchroniczny: działa tylko na zboczu
  else if (en) q <= d;       // en = 0: q <= q
end
\end{sv}

\begin{figure}[h]
\centering
\begin{tikzpicture}
  \pic (m1) at (0,0) {muks};
  \pic (m2) at (2.2,0) {muks};
  \node[draw, thick, fill=zadaniejasny, minimum width=11mm, minimum height=14mm, font=\sffamily\scriptsize] (ff) at (4.6,0) {D\quad Q};
  \draw[<-, thick] (m1-in1) -- ++(-8mm,0) node[left, sygnal]{d};
  \draw[<-, thick] (m1-sel) -- ++(0,-6mm) node[below, sygnal]{en};
  \draw[->, thick] (m1-out) -- ++(5mm,0) |- (m2-in0);
  \draw[<-, thick] (m2-in1) -- ++(-5mm,0) node[left, sygnal]{0};
  \draw[<-, thick] (m2-sel) -- ++(0,-6mm) node[below, sygnal]{rst};
  \draw[->, thick] (m2-out) -- (ff.west);
  \draw[->, thick] (ff.east) -- ++(10mm,0) node[right, sygnal]{q};
  \draw[thick] ($(ff.east)+(4mm,0)$) -- ++(0,-19mm) -| ($(m1-in0)+(-6mm,0)$) -- (m1-in0);
  \draw[<-, thick] (ff.north) -- ++(0,6mm) node[above, sygnal]{clk};
\end{tikzpicture}
\caption{Rejestr z resetem synchronicznym i zezwoleniem: dwa multipleksery i przerzutnik D.}
\label{fig:regen}
\end{figure}

Sprzętowo (rys.~\ref{fig:regen}) to przerzutnik D z dwoma multiplekserami
przed nim. Biblioteki FPGA i ASIC mają gotowe przerzutniki z wejściem
\emph{enable} (\code{\$\_SDFFE\_} z rozdziału~\ref{ch:start}).

\paragraph{Reset synchroniczny czy asynchroniczny?} Asynchroniczny
(\code{always\_ff @(posedge clk or posedge rst)}) działa bez zegara, ale
komplikuje analizę czasową przy zdejmowaniu resetu. W kursie używamy tylko
synchronicznego, który na FPGA jest standardem.

\paragraph{Wartości początkowe.} Na FPGA przerzutniki dostają wartości z
bitstreamu, więc \code{initial} i \code{logic q = 0;} są syntezowalne. W ASIC
po włączeniu zasilania stan jest losowy i tylko reset daje pewność. Dlatego
wszystko, co musi mieć znaną wartość (PC, stan automatu, bity \code{valid}),
resetujemy jawnie.

\section{Czas w układzie synchronicznym}

\begin{figure}[h]
\centering
\begin{tikzpicture}
  \node[draw, thick, fill=zadaniejasny, minimum width=10mm, minimum height=12mm, font=\sffamily\scriptsize] (f1) at (0,0) {FF1};
  \node[blok, minimum width=34mm, minimum height=12mm] (c) at (4,0) {logika\\kombinacyjna};
  \node[draw, thick, fill=zadaniejasny, minimum width=10mm, minimum height=12mm, font=\sffamily\scriptsize] (f2) at (8,0) {FF2};
  \draw[->, thick] (f1) -- node[above, etykieta]{$t_{clk\to q}$} (c);
  \draw[->, thick] (c) -- node[above, etykieta]{$t_{comb}$} (f2);
  \draw[thick] (0,-1.5) -- (8,-1.5);
  \draw[->, thick] (0,-1.5) -- (f1.south);
  \draw[->, thick] (8,-1.5) -- (f2.south);
  \node[sygnal, below] at (4,-1.5) {clk};
\end{tikzpicture}
\caption{Ścieżka między dwoma rejestrami.}
\label{fig:timing}
\end{figure}

Żeby FF2 poprawnie zapamiętał wynik, sygnał musi do niego dotrzeć przed
zboczem z zapasem $t_{setup}$:
\[ T_{clk} \;\geq\; t_{clk\to q} + t_{comb}^{max} + t_{setup}. \]

Najdłuższa ścieżka kombinacyjna między dwoma rejestrami (\textbf{ścieżka
krytyczna}) wyznacza maksymalną częstotliwość zegara. Tak wygląda
„złożoność obliczeniowa” sprzętu: nie liczba operacji, tylko \emph{głębokość
logiki między rejestrami}. Za długą ścieżkę można przeciąć rejestrem. Na tym
polega \textbf{potokowanie} (rozdział~\ref{ch:pipe}).

W symulacji RTL czasów nie ma, bo logika „liczy się w zerowym czasie”. O
timing dbają synteza i place\&route. Przybliżenie pokazuje \code{make synth}
(\emph{longest topological path}).

\section{Pamięci}

\begin{sv}
logic [31:0] mem [0:255];                       // 256 słów po 32 bity
always_ff @(posedge clk)
  if (we) mem[waddr] <= wdata;                  // zapis synchroniczny
assign rdata = mem[raddr];                      // odczyt asynchroniczny
// albo:
always_ff @(posedge clk) rdata <= mem[raddr];   // odczyt synchroniczny
\end{sv}

\begin{description}
  \item[Odczyt synchroniczny] (dana dostępna takt później): tak działają bloki
    RAM w FPGA (BRAM) i pamięci SRAM. Są gęste i tanie.
  \item[Odczyt asynchroniczny] wymaga zbudowania pamięci z przerzutników i
    multiplekserów albo z „rozproszonej” pamięci LUT. Dla dużych pamięci jest
    drogi.
\end{description}

Pamięć w testbenchu naszego procesora ma odczyt asynchroniczny, bo upraszcza
to procesor jednocyklowy. Bank rejestrów (32×32 bity) jest mały i też ma
odczyt asynchroniczny. W rozdziale 3 zobaczysz, ile to kosztuje.

\section{Automaty skończone (FSM)}

Automat to rejestr stanu i dwie funkcje kombinacyjne:
\emph{następny stan} = $f$(stan, wejścia) oraz \emph{wyjścia} = $g$(stan)
w automacie \textbf{Moore'a} albo $g$(stan, wejścia) w automacie
\textbf{Mealy'ego}. Wyjścia Moore'a zmieniają się tylko po zboczu zegara (są
„czyste”, ale opóźnione o takt). Mealy reaguje w tym samym takcie, ale jego
wyjście zależy kombinacyjnie od wejść (dłuższe ścieżki, ryzyko pętli
kombinacyjnych między modułami).

\subsection*{Przykład: detektor sekwencji 110 (Moore, z nakładaniem)}

\begin{figure}[h]
\centering
\begin{tikzpicture}[->, >=Stealth, node distance=27mm, thick, every state/.style={fill=akcentjasny, minimum size=13mm, font=\sffamily\scriptsize, align=center}]
  \node[state, initial, initial text=rst] (s0) {S0\\$\varnothing$};
  \node[state, right of=s0] (s1) {S1\\„1”};
  \node[state, right of=s1] (s11) {S11\\„11”};
  \node[state, right of=s11, fill=zadaniejasny, double] (s110) {S110\\found=1};
  \path (s0) edge[loop above] node[sygnal]{0} ()
             edge node[above, sygnal]{1} (s1)
        (s1) edge node[above, sygnal]{1} (s11)
             edge[bend left=35] node[below, sygnal]{0} (s0)
        (s11) edge[loop above] node[sygnal]{1} ()
              edge node[above, sygnal]{0} (s110)
        (s110) edge[bend right=45] node[above, sygnal]{1} (s1)
               edge[bend left=55] node[below, sygnal]{0} (s0);
\end{tikzpicture}
\caption{Automat Moore'a wykrywający 110. Wyjście \code{found} zależy tylko od stanu.}
\label{fig:fsm110}
\end{figure}

Każda strzałka na diagramie to jedna gałąź w \code{case}. Styl obowiązujący w
kursie, czyli \textbf{dwa procesy}:

\begin{sv}
typedef enum logic [1:0] {S0, S1, S11, S110} state_t;
state_t state, next;

always_ff @(posedge clk)              // 1) tylko rejestr stanu
  if (rst) state <= S0;
  else     state <= next;

always_comb begin                     // 2) tylko logika następnego stanu
  next = state;                       //    wartość domyślna (brak zatrzasku!)
  case (state)
    S0:   if (in) next = S1;   else next = S0;
    S1:   if (in) next = S11;  else next = S0;
    S11:  if (in) next = S11;  else next = S110;
    S110: if (in) next = S1;   else next = S0;
    default: next = S0;
  endcase
end

assign found = (state == S110);       // wyjście Moore'a
\end{sv}

\begin{pulapka}[Icarus i \texttt{enum}]
\code{next = in ? S1 : S0;} daje w Icarusie błąd \emph{This assignment requires
an explicit cast}. Pisz \code{if/else} (jak wyżej) albo
\code{next = state\_t'(in ? S1 : S0);}.
\end{pulapka}

\subsection*{Automat z licznikiem: sterowanie i ścieżka danych}

Większość prawdziwych układów (UART, kontroler pamięci, procesor wielocyklowy)
to mały automat sterujący i „ścieżka danych”: liczniki i rejestry przesuwne.
Automat mówi „licz”, „przesuń”, „załaduj”, a ścieżka danych odpowiada
„doliczyłem do końca”. Nadajnik UART z zadania 2.5 jest dokładnie taki.

\begin{figure}[h]
\centering
\begin{tikzpicture}[x=5.2mm, y=6mm]
  \node[sygnal, anchor=east] at (-0.4,0.25) {tx};
  \draw[thick] (-1,0.5) -- (0,0.5) -- (0,0) -- (2,0);
  \foreach \i/\b in {0/1,1/0,2/1,3/1,4/0,5/0,6/1,7/0} {
    \draw[thick] ({2+2*\i},0) -- ({2+2*\i},0.5) -- ({4+2*\i},0.5) -- ({4+2*\i},0) -- cycle;
    \node[sygnal] at ({3+2*\i},0.25) {d\i};
  }
  \draw[thick] (18,0.5) -- (22,0.5);
  \node[etykieta] at (1,-0.5) {start=0};
  \node[etykieta] at (19,-0.5) {stop=1};
  \draw[decorate, decoration={brace, amplitude=4pt}] (2,0.8) -- node[above=4pt, etykieta]{8 bitów danych, najmłodszy pierwszy} (18,0.8);
  \draw[decorate, decoration={brace, amplitude=3pt, mirror}] (0,-1) -- node[below=3pt, etykieta]{CLKS\_PER\_BIT taktów} (2,-1);
\end{tikzpicture}
\caption{Ramka UART 8N1. Linia w spoczynku ma stan 1.}
\label{fig:uart}
\end{figure}

\section{Porównania z parametrami}

\code{parameter int DEPTH = 8} ma 32 bity. Porównanie \code{count == DEPTH},
gdzie \code{count} ma 4 bity, jest poprawne, ale Verilator zgłosi
\code{WIDTHEXPAND}. Czysty idiom to stała o właściwej szerokości:

\begin{sv}
localparam int AW = $clog2(DEPTH);
localparam logic [AW:0] FULL = (AW+1)'(DEPTH);   // rzutowanie na AW+1 bitów
assign full = (count == FULL);
\end{sv}

\section{Kolejka FIFO: wskaźniki z dodatkowym bitem}

FIFO o głębokości $D = 2^k$ to pamięć i dwa wskaźniki: zapisu i odczytu.
Gdyby wskaźniki miały $k$ bitów, stan „pełna” i „pusta” wyglądałyby tak samo
(wskaźniki równe). Sztuczka: wskaźniki mają $k+1$ bitów, a pamięć adresuje się
młodszymi $k$ bitami. Wtedy:

\begin{itemize}
  \item \code{count = wr\_ptr - rd\_ptr}: odejmowanie modulo $2^{k+1}$ zawija się samo,
  \item \code{empty}: wskaźniki równe,
  \item \code{full}: różnią się tylko najstarszym bitem (czyli \code{count == D}).
\end{itemize}

\section{Sygnały z zewnątrz (na przyszłość)}

Przycisk albo linia RX z UART zmieniają się niezależnie od Twojego zegara.
Przerzutnik próbkujący taki sygnał może wpaść w \textbf{metastabilność}.
Standardowe lekarstwo to dwa przerzutniki pod rząd (synchronizator) przed
jakąkolwiek logiką. W symulacji tego nie zobaczysz, a na płytce objawia się
„losowymi” błędami raz na godzinę.

\section{Zadania}

\begin{zadanie}{2.1 --- Detektor zbocza (\code{rtl/edge\_detect.sv})}
Jeden przerzutnik i jedna bramka. W przebiegach (\code{make wave T=edge\_detect})
zobacz, że \code{rise} jest kombinacyjny: pojawia się razem ze zmianą
\code{in}, a znika na zboczu zegara.
\end{zadanie}

\begin{zadanie}{2.2 --- LFSR (\code{rtl/lfsr16.sv})}
Wielomian $x^{16}+x^{14}+x^{13}+x^{11}+1$, ziarno \code{16'hACE1}. Test
sprawdza też okres (65535). Dlaczego okres to $2^{16}-1$, a nie $2^{16}$? Jakiego
stanu LFSR nigdy nie osiągnie i co by się stało, gdyby w nim wystartował?
\end{zadanie}

\begin{zadanie}{2.3 --- Detektor 1011 (\code{rtl/seq\_detect.sv})}
Automat Moore'a, 5 stanów, styl dwuprocesowy. \textbf{Najpierw narysuj
diagram} jak na rys.~\ref{fig:fsm110}. Najtrudniejsze są przejścia przy
nakładaniu: dokąd przejść ze stanu „1011”, gdy przyjdzie 0, a dokąd, gdy 1?
\end{zadanie}

\begin{zadanie}{2.4 --- FIFO (\code{rtl/fifo.sv})}
Test porównuje FIFO z kolejką \code{q[\$]} z SystemVeriloga przez 3000 losowych
taktów, także przy jednoczesnym \code{push} i \code{pop}.
\end{zadanie}

\begin{zadanie}{2.5 --- Nadajnik UART (\code{rtl/uart\_tx.sv})}
Ramka z rys.~\ref{fig:uart}. Test to „odbiornik programowy” w testbenchu (warto
przeczytać: \code{fork/join}, taski czekające na zbocza). Po zrobieniu:
\code{make synth T=uart\_tx}. Czy umiesz przypisać każdy przerzutnik do
konkretnego rejestru w swoim kodzie? Ten moduł przyda się na FPGA (rozdział 9).
\end{zadanie}

\begin{gotowe}
\code{make test} → 5 × PASS, \code{make lint} czysto.
\end{gotowe}

\begin{kontrolne}
  \item Co dokładnie zrobi symulator z \code{a <= b; b <= a;} na zboczu zegara? A z \code{a = b; b = a;}?
  \item Od czego zależy maksymalna częstotliwość zegara układu? Co to jest ścieżka krytyczna?
  \item Czym różni się odczyt synchroniczny od asynchronicznego i który z nich mają bloki RAM w FPGA?
  \item Narysuj automat Mealy'ego wykrywający 110. Ile ma stanów? Kiedy zmienia się jego wyjście?
  \item Dlaczego wskaźniki FIFO mają o jeden bit więcej, niż potrzeba do adresowania?
\end{kontrolne}

\chapter{ALU, bank rejestrów, komparator skoków}
\label{ch:alu}

\metryczka{zbudować trzy największe bloki ścieżki danych RV32I. Od tego rozdziału
interfejsy modułów są kontraktem: procesory z rozdziałów 5 i 8 użyją
dokładnie tych portów.}{1--2 wieczory}{03-alu-regfile/}

\section{RV32I w pigułce}

RISC-V RV32I to 32-bitowa architektura typu \emph{load/store}:
\begin{itemize}
  \item \textbf{32 rejestry} \code{x0}--\code{x31} po 32 bity, przy czym
    \textbf{\code{x0} jest zawsze zerem}, a zapis do niego jest ignorowany.
    Dzięki temu wiele pseudoinstrukcji nic nie kosztuje:
    \code{mv a,b} = \code{addi a,b,0}, \code{nop} = \code{addi x0,x0,0},
    \code{j L} = \code{jal x0,L}.
  \item Arytmetyka działa tylko na rejestrach, a do pamięci dostęp jest
    wyłącznie przez load i store.
  \item \textbf{Brak flag}: nie ma rejestru statusu jak w x86 czy ARM. Skoki
    warunkowe same porównują dwa rejestry (\code{blt a0, a1, L}), dlatego
    zamiast flag z ALU budujemy osobny komparator.
\end{itemize}

\begin{center}
\small
\begin{tabular}{@{}lll@{}}
\toprule
\textbf{Rejestr} & \textbf{Nazwa ABI} & \textbf{Rola} \\
\midrule
x0 & zero & stałe zero \\
x1 & ra & adres powrotu \\
x2 & sp & wskaźnik stosu \\
x3 & gp & global pointer (w testach kursu: numer podtestu) \\
x5--7, x28--31 & t0--t6 & tymczasowe (\emph{caller-saved}) \\
x8--9, x18--27 & s0--s11 & zachowywane (\emph{callee-saved}), s0 = fp \\
x10--17 & a0--a7 & argumenty i wynik funkcji \\
\bottomrule
\end{tabular}
\end{center}

\section{ALU}

Dziesięć operacji wystarcza na całą arytmetykę RV32I:

\begin{center}
\small
\begin{tabularx}{\linewidth}{@{}l l X@{}}
\toprule
\textbf{Kod \code{op}} & \textbf{Operacja} & \textbf{Używana przez} \\
\midrule
\code{ALU\_ADD} 0\_000 & $a + b$ & add, addi, adresy load/store, auipc, lui ($0+imm$), jalr \\
\code{ALU\_SUB} 1\_000 & $a - b$ & sub \\
\code{ALU\_SLL} 0\_001 & $a \ll b[4{:}0]$ & sll, slli \\
\code{ALU\_SLT} 0\_010 & $a <_s b$ ? 1 : 0 & slt, slti \\
\code{ALU\_SLTU} 0\_011 & $a <_u b$ ? 1 : 0 & sltu, sltiu \\
\code{ALU\_XOR} 0\_100 & $a \oplus b$ & xor, xori \\
\code{ALU\_SRL} 0\_101 & $a \gg b[4{:}0]$ (zera) & srl, srli \\
\code{ALU\_SRA} 1\_101 & $a \gg b[4{:}0]$ (znak) & sra, srai \\
\code{ALU\_OR} 0\_110 & $a \lor b$ & or, ori \\
\code{ALU\_AND} 0\_111 & $a \land b$ & and, andi \\
\bottomrule
\end{tabularx}
\end{center}

Kody są celowo równe \code{\{funct7[5], funct3\}} instrukcji R-type. Dekoder z
rozdziału 5 dla \code{add/sub/sll/...} po prostu przepisze te bity z instrukcji.
Stałe pochodzą z pakietu, który importujesz w nagłówku modułu:

\begin{sv}
module alu import rv32i_pkg::*; (
  input  logic [31:0] a, b,
  input  logic [3:0]  op,
  output logic [31:0] y
);
  always_comb begin
    case (op)
      ALU_ADD: y = a + b;
      // ...
\end{sv}

\begin{pulapka}[Trzy klasyczne błędy w ALU]
(1) \code{a >>> n} bez \code{\$signed(a)} to zwykłe przesunięcie logiczne.
(2) \code{a < b} na \code{logic} porównuje bez znaku, więc SLT wymaga
\code{\$signed}. (3) Przesuwa się o \code{b[4:0]}, nie o całe \code{b}: ISA
definiuje \code{sll x1,x2,x3} jako przesunięcie o \code{x3 \& 31}.
\end{pulapka}

\subsection*{Współdzielony sumator (dla chętnych)}

Prosta wersja z \code{case} buduje osobny sumator dla ADD, osobny dla SUB i
jeszcze jeden dla SLT/SLTU. Wszystkie da się obsłużyć \emph{jednym}:
$a - b = a + \overline{b} + 1$, SLTU to brak przeniesienia z odejmowania, a SLT
to znak różnicy poprawiony o przepełnienie:
\[ \mathrm{SLT} = (a_{31} \ne b_{31}) \;?\; a_{31} \;:\; (a-b)_{31}. \]
Porównaj \code{make synth T=alu} przed i po. Yosys część takiego
współdzielenia robi sam, części nie.

\section{Bank rejestrów}

\begin{figure}[h]
\centering
\begin{tikzpicture}
  \node[draw, thick, fill=zadaniejasny, minimum width=26mm, minimum height=40mm, align=center, font=\sffamily\small] (rf) {x1 \\ x2 \\ \vdots \\ x31};
  \node[blok, minimum width=16mm, minimum height=12mm] (r1) at ($(rf.east)+(30mm,14mm)$) {mux\\32:1};
  \node[blok, minimum width=16mm, minimum height=12mm] (r2) at ($(rf.east)+(30mm,-14mm)$) {mux\\32:1};
  \draw[->, thick] ($(rf.east)+(0,10mm)$) -- ++(8mm,0) |- ($(r1.west)+(0,0)$);
  \draw[->, thick] ($(rf.east)+(0,-10mm)$) -- ++(8mm,0) |- ($(r2.west)+(0,0)$);
  \draw[->, thick] (r1.east) -- ++(10mm,0) node[right, sygnal]{rs1\_val};
  \draw[->, thick] (r2.east) -- ++(10mm,0) node[right, sygnal]{rs2\_val};
  \draw[<-, thick] (r1.north) -- ++(0,6mm) node[above, sygnal]{rs1};
  \draw[<-, thick] (r2.south) -- ++(0,-6mm) node[below, sygnal]{rs2};
  \node[blok, left=14mm of rf, minimum width=20mm] (dec) {dekoder\\zapisu};
  \draw[->, thick] (dec) -- node[above, etykieta]{we[i]} (rf);
  \draw[<-, thick] (dec.north) -- ++(0,6mm) node[above, sygnal]{we, rd};
  \draw[<-, thick] (rf.south) -- ++(0,-8mm) node[below, sygnal]{rd\_val, clk};
\end{tikzpicture}
\caption{Bank rejestrów: 31 rejestrów (x0 jest stałą), dwa multipleksery
odczytu 32:1 i dekoder zapisu.}
\label{fig:rf}
\end{figure}

\begin{itemize}
  \item \textbf{x0}: najprościej trzymać tablicę \code{regs[1:31]}, przy
    odczycie dla \code{rs == 0} zwracać 0, a zapis z \code{rd == 0} ignorować.
  \item \textbf{Odczyt w takcie zapisu} pod ten sam adres zwraca \emph{starą}
    wartość. W procesorze jednocyklowym to poprawne: instrukcja czyta rejestry
    na początku taktu, a wynik zapisuje na jego końcu. W potoku (rozdział 8)
    „bypass” dodasz obok banku, nie wewnątrz.
  \item \textbf{Zerowanie w \code{initial}}. Prawdziwy procesor ma po resecie
    śmieci w rejestrach. U nas zerujemy je z powodu praktycznego: w rozdziale 6
    porównujemy ślad procesora z emulatorem, który startuje od zer. Kompilatory
    C zapisują na stos niezainicjowane rejestry \code{s*} w prologach funkcji,
    więc RTL miałby tam \code{X}, emulator 0, i porównanie by się rozjechało.
    Na FPGA \code{initial} jest syntezowalne, w ASIC nie.
\end{itemize}

\subsection*{Ile to kosztuje?}

Synteza wzorcowego rozwiązania daje 992 przerzutniki (31×32) i ok. 1900
multiplekserów 2:1. Każdy port odczytu to 32 drzewa po 31 multiplekserów.
Bank rejestrów jest \emph{dwa razy większy niż cała ALU}! Dlatego prawdziwe
procesory na FPGA często trzymają rejestry w BRAM-ie z odczytem
synchronicznym, a odczyt rejestrów jest wtedy osobnym etapem potoku.

\section{Komparator skoków}

\begin{center}
\small
\begin{tabular}{@{}llll@{}}
\toprule
\code{funct3} & Instrukcja & Warunek skoku & Uwagi \\
\midrule
000 & beq & $a = b$ & \\
001 & bne & $a \ne b$ & negacja beq \\
100 & blt & $a <_s b$ & ze znakiem \\
101 & bge & $a \geq_s b$ & negacja blt \\
110 & bltu & $a <_u b$ & bez znaku \\
111 & bgeu & $a \geq_u b$ & negacja bltu \\
\bottomrule
\end{tabular}
\end{center}

Struktura kodowania: \code{funct3[0]} neguje wynik, a \code{funct3[2:1]}
wybiera rodzaj porównania. Da się to zapisać krócej niż przez sześć przypadków
\code{case}.

\section{Zadania}

\begin{zadanie}{3.1 --- ALU (\code{rtl/alu.sv})}
Test: 20\,000 kombinacji z naciskiem na przypadki brzegowe (0, 1, $-1$,
\code{0x80000000}, \code{0x7FFFFFFF}, przesunięcie o 33).
\end{zadanie}

\begin{zadanie}{3.2 --- Bank rejestrów (\code{rtl/regfile.sv})}
Test sprawdza zerowy stan początkowy, x0, brak zapisu przy \code{we = 0} i
\emph{stary} odczyt w takcie zapisu.
\end{zadanie}

\begin{zadanie}{3.3 --- Komparator (\code{rtl/branch\_cmp.sv})}
\end{zadanie}

\begin{zadanie}{3.4 --- Synteza i koszt}
\code{make synth T=alu}, \code{T=regfile}, \code{T=branch\_cmp}. Zapisz liczby
komórek i długości ścieżek. Który blok jest największy i dlaczego? Która
operacja ALU wyznacza najdłuższą ścieżkę? (Zakomentuj na chwilę ADD/SUB/SLT
i zsyntezuj ponownie.)
\end{zadanie}

\begin{gotowe}
\code{make test} → 3 × PASS, \code{make lint} czysto.
\end{gotowe}

\begin{kontrolne}
  \item Dlaczego RISC-V nie ma flag i jakie to ma konsekwencje dla sprzętu?
  \item Jaki wynik da \code{SLT} dla $a = \code{0x80000000}$, $b = 1$? A \code{SLTU}?
  \item Dlaczego odczyt w takcie zapisu ma zwracać starą wartość w procesorze jednocyklowym?
  \item Ile multiplekserów 2:1 potrzeba na jeden 32-bitowy port odczytu z 32 rejestrów?
\end{kontrolne}

\chapter{Architektura RV32I i emulator w C}
\label{ch:isa}

\metryczka{znać RV32I na poziomie bitów i napisać \emph{model referencyjny}
procesora w C: emulator, z którym w rozdziałach 6--8 porównasz swój sprzęt
instrukcja po instrukcji.}{2--3 wieczory}{04-isa-emulator/}

\section{Formaty instrukcji}

Każda instrukcja ma 32 bity. \textbf{Pola rejestrów są zawsze w tych samych
miejscach}, więc dekoder może czytać rejestry, zanim jeszcze rozpozna
instrukcję.

% \pole{y}{starszy bit}{młodszy bit}{etykieta}{kolor}
\newcommand{\pole}[5]{%
  \draw[thick, fill=#5] ({(31-#2)*0.42},#1) rectangle ({(32-#3)*0.42},{#1+0.62});
  \node[font=\sffamily\scriptsize] at ({(31-#2+32-#3)*0.21},{#1+0.31}) {#4};}

\begin{figure}[h]
\centering
\begin{tikzpicture}
  \foreach \b in {31,25,24,20,19,15,14,12,11,7,6,0} {
    \node[font=\ttfamily\tiny, text=szary] at ({(31-\b+0.5)*0.42},0.85) {\b};
  }
  % R
  \pole{0}{31}{25}{funct7}{akcentjasny} \pole{0}{24}{20}{rs2}{zadaniejasny}
  \pole{0}{19}{15}{rs1}{zadaniejasny} \pole{0}{14}{12}{f3}{akcentjasny}
  \pole{0}{11}{7}{rd}{zadaniejasny} \pole{0}{6}{0}{opcode}{akcentjasny}
  \node[anchor=west, font=\sffamily\small\bfseries] at (13.7,0.31) {R};
  % I
  \pole{-0.9}{31}{20}{imm[11:0]}{pulapkajasny} \pole{-0.9}{19}{15}{rs1}{zadaniejasny}
  \pole{-0.9}{14}{12}{f3}{akcentjasny} \pole{-0.9}{11}{7}{rd}{zadaniejasny}
  \pole{-0.9}{6}{0}{opcode}{akcentjasny}
  \node[anchor=west, font=\sffamily\small\bfseries] at (13.7,-0.59) {I};
  % S
  \pole{-1.8}{31}{25}{imm[11:5]}{pulapkajasny} \pole{-1.8}{24}{20}{rs2}{zadaniejasny}
  \pole{-1.8}{19}{15}{rs1}{zadaniejasny} \pole{-1.8}{14}{12}{f3}{akcentjasny}
  \pole{-1.8}{11}{7}{imm[4:0]}{pulapkajasny} \pole{-1.8}{6}{0}{opcode}{akcentjasny}
  \node[anchor=west, font=\sffamily\small\bfseries] at (13.7,-1.49) {S};
  % B
  \pole{-2.7}{31}{31}{\tiny 12}{pulapkajasny} \pole{-2.7}{30}{25}{imm[10:5]}{pulapkajasny}
  \pole{-2.7}{24}{20}{rs2}{zadaniejasny} \pole{-2.7}{19}{15}{rs1}{zadaniejasny}
  \pole{-2.7}{14}{12}{f3}{akcentjasny} \pole{-2.7}{11}{8}{imm[4:1]}{pulapkajasny}
  \pole{-2.7}{7}{7}{\tiny 11}{pulapkajasny} \pole{-2.7}{6}{0}{opcode}{akcentjasny}
  \node[anchor=west, font=\sffamily\small\bfseries] at (13.7,-2.39) {B};
  % U
  \pole{-3.6}{31}{12}{imm[31:12]}{pulapkajasny} \pole{-3.6}{11}{7}{rd}{zadaniejasny}
  \pole{-3.6}{6}{0}{opcode}{akcentjasny}
  \node[anchor=west, font=\sffamily\small\bfseries] at (13.7,-3.29) {U};
  % J
  \pole{-4.5}{31}{31}{\tiny 20}{pulapkajasny} \pole{-4.5}{30}{21}{imm[10:1]}{pulapkajasny}
  \pole{-4.5}{20}{20}{\tiny 11}{pulapkajasny} \pole{-4.5}{19}{12}{imm[19:12]}{pulapkajasny}
  \pole{-4.5}{11}{7}{rd}{zadaniejasny} \pole{-4.5}{6}{0}{opcode}{akcentjasny}
  \node[anchor=west, font=\sffamily\small\bfseries] at (13.7,-4.19) {J};
\end{tikzpicture}
\caption{Sześć formatów instrukcji RV32I (f3 = funct3). Na czerwono fragmenty stałej natychmiastowej.}
\label{fig:formaty}
\end{figure}

\section{Stałe natychmiastowe (immediate)}

\begin{center}
\small
\begin{tabular}{@{}ll@{}}
\toprule
Format & Immediate jako 32-bitowa liczba (\code{i} = instrukcja) \\
\midrule
I & \code{\{\{20\{i[31]\}\}, i[31:20]\}} \\
S & \code{\{\{20\{i[31]\}\}, i[31:25], i[11:7]\}} \\
B & \code{\{\{19\{i[31]\}\}, i[31], i[7], i[30:25], i[11:8], 1'b0\}} \\
U & \code{\{i[31:12], 12'b0\}} \\
J & \code{\{\{11\{i[31]\}\}, i[31], i[19:12], i[20], i[30:21], 1'b0\}} \\
\bottomrule
\end{tabular}
\end{center}

\paragraph{Dlaczego bity są tak „pomieszane”?} Zasada projektowa RISC-V: każdy
bit wyniku ma pochodzić z jak najmniejszej liczby pozycji w instrukcji.
\begin{itemize}
  \item Bit znaku to \emph{zawsze} \code{instr[31]}, więc rozszerzanie znakiem
    zaczyna się, zanim dekoder rozpozna format.
  \item Bity \code{imm[10:5]} są zawsze w \code{instr[30:25]}, a \code{imm[4:1]}
    w jednym z dwóch miejsc.
  \item W sprzęcie każdy bit immediate'a to mały multiplekser o 2--3 wejściach
    zamiast dużego.
\end{itemize}
W formatach B i J nie ma bitu 0, bo skoki są zawsze o parzystą liczbę bajtów
(rozszerzenie C ma instrukcje 16-bitowe). Dzięki temu zasięg jest dwa razy
większy: $\pm$4 KiB dla B i $\pm$1 MiB dla J.

\section{Lista instrukcji RV32I}

\begin{center}
\small
\begin{tabularx}{\linewidth}{@{}l c l X@{}}
\toprule
\textbf{opcode} & \textbf{Format} & \textbf{Instrukcje} & \textbf{Semantyka} \\
\midrule
\code{0110111} LUI & U & lui rd, imm & rd = imm $\ll$ 12 \\
\code{0010111} AUIPC & U & auipc rd, imm & rd = pc + (imm $\ll$ 12) \\
\code{1101111} JAL & J & jal rd, off & rd = pc+4; pc += off \\
\code{1100111} JALR & I & jalr rd, off(rs1) & t = (rs1+off) \& \textasciitilde1; rd = pc+4; pc = t \\
\code{1100011} BRANCH & B & beq bne blt bge bltu bgeu & jeśli warunek(rs1, rs2): pc += off \\
\code{0000011} LOAD & I & lb lh lw lbu lhu & rd = rozszerz(mem[rs1+off]) \\
\code{0100011} STORE & S & sb sh sw & mem[rs1+off] = rs2 (8/16/32 b) \\
\code{0010011} OP-IMM & I & addi slti sltiu xori ori andi slli srli srai & rd = rs1 op imm \\
\code{0110011} OP & R & add sub sll slt sltu xor srl sra or and & rd = rs1 op rs2 \\
\code{0001111} FENCE & I & fence & u nas: nop \\
\code{1110011} SYSTEM & I & ecall, ebreak, csr* & u nas: nop \\
\bottomrule
\end{tabularx}
\end{center}

\code{funct3} wybiera wariant: rodzaj operacji ALU, szerokość load/store albo
warunek skoku. \code{funct7 = 0100000} (czyli \code{instr[30] = 1}) odróżnia
\code{sub} od \code{add} oraz \code{sra/srai} od \code{srl/srli}. Przy
\code{slli/srli/srai} pole \code{rs2} zawiera \code{shamt}.

\begin{pulapka}[\texttt{jalr} z \texttt{rd == rs1}]
W \code{jalr a0, 0(a0)} cel liczysz ze \emph{starej} wartości \code{a0}, a
dopiero potem zapisujesz do niej adres powrotu. W C łatwo napisać to w złej
kolejności. Test \code{06\_jal\_jalr} to sprawdza.
\end{pulapka}

\section{Pseudoinstrukcje i ładowanie stałych}

\begin{center}
\small
\begin{tabular}{@{}ll@{\hspace{2em}}ll@{}}
\toprule
\textbf{Pseudo} & \textbf{Rozwinięcie} & \textbf{Pseudo} & \textbf{Rozwinięcie} \\
\midrule
\code{nop} & \code{addi x0, x0, 0} & \code{j L} & \code{jal x0, L} \\
\code{mv rd, rs} & \code{addi rd, rs, 0} & \code{call L} & \code{jal ra, L} \\
\code{not rd, rs} & \code{xori rd, rs, -1} & \code{ret} & \code{jalr x0, 0(ra)} \\
\code{neg rd, rs} & \code{sub rd, x0, rs} & \code{beqz rs, L} & \code{beq rs, x0, L} \\
\code{seqz rd, rs} & \code{sltiu rd, rs, 1} & \code{bgt a, b, L} & \code{blt b, a, L} \\
\bottomrule
\end{tabular}
\end{center}

\code{li rd, imm32} rozwija się w \code{lui rd, hi} + \code{addi rd, rd, lo}.
\textbf{Pułapka:} \code{addi} rozszerza 12-bitową stałą \emph{znakiem}, więc dla
\code{0x12345FFF} nie wystarczy \code{lui 0x12345; addi 0xFFF}, bo
\code{0xFFF} to $-1$. Poprawnie: $hi = (v + \code{0x800}) \gg 12$, czyli
\code{lui 0x12346; addi -1}.

\section{Pamięć}

Adresowanie jest bajtowe, \textbf{little-endian}: słowo \code{0x11223344} pod
adresem 0 to bajty \code{44 33 22 11} pod adresami 0, 1, 2, 3.
\code{lb}/\code{lh} rozszerzają znakiem, a \code{lbu}/\code{lhu} zerami.
Zakładamy dostępy \emph{wyrównane}.

\begin{figure}[h]
\centering
\begin{tikzpicture}[x=13mm]
  \foreach \i/\b in {0/44,1/33,2/22,3/11} {
    \draw[thick, fill=akcentjasny] (\i,0) rectangle ++(1,0.7);
    \node[font=\ttfamily\small] at (\i+0.5,0.35) {\b};
    \node[sygnal, text=szary] at (\i+0.5,-0.25) {adr \i};
  }
  \draw[thick, fill=zadaniejasny] (6,0) rectangle ++(4,0.7);
  \foreach \i/\b in {0/11,1/22,2/33,3/44} { \node[font=\ttfamily\small] at (6.5+\i,0.35) {\b}; }
  \node[sygnal, text=szary] at (6.5,-0.25) {[31:24]};
  \node[sygnal, text=szary] at (9.5,-0.25) {[7:0]};
  \draw[->, thick] (4.2,0.35) -- node[above, etykieta]{lw} (5.8,0.35);
\end{tikzpicture}
\caption{Little-endian: najmłodszy bajt słowa pod najniższym adresem.}
\end{figure}

\section{Nasza maszyna}

To samo „środowisko” symulują emulator (\plik{emu/rv32sim.c}) i testbench RTL
(\plik{common/tb/soc\_tb.sv}):

\begin{center}
\small
\begin{tabularx}{\linewidth}{@{}l X l@{}}
\toprule
\textbf{Adres} & \textbf{Co} & \textbf{Dostęp} \\
\midrule
\code{0x0000\_0000}--\code{0x0000\_FFFF} & RAM 64 KiB, tu ładowany jest program & R/W \\
\code{0x1000\_0000} & TOHOST: zapis kończy symulację; 1 = PASS, $(n \ll 1)|1$ = FAIL w podteście $n$ & W \\
\code{0x1000\_0004} & PUTCHAR: wypisuje znak \code{wdata[7:0]} & W \\
\code{0x1000\_0008} & CYCLES: licznik cykli (w emulatorze: liczba instrukcji) & R \\
\bottomrule
\end{tabularx}
\end{center}

Po resecie \code{pc = 0} i wszystkie rejestry mają wartość 0. Testy
(\plik{common/sw/tests/*.S}) trzymają w rejestrze \code{gp} numer podtestu i
kończą skokiem do \code{pass} albo \code{fail} (\plik{test\_end.S}):

\begin{asm}
    li   gp, 3              # podtest 3: slti / sltiu
    slti a1, a0, 6
    li   t2, 1
    bne  a1, t2, fail       # przy błędzie: TOHOST = (3 << 1) | 1
    ...
    j    pass
.include "test_end.S"
\end{asm}

\section{Asembler i disasembler kursu}

\begin{sh}
T=../common/tools/rvtool.py
python3 $T asm ../common/sw/tests/08_fib.S -I ../common/sw/tests -o fib.hex -l fib.lst
python3 $T dis fib.hex              # disasemblacja całego obrazu
python3 $T dis 0x00a50533           # jedno słowo: add a0, a0, a0
\end{sh}

Format \code{.hex}: jedno 32-bitowe słowo na linię, od adresu 0 (to, co czyta
\code{\$readmemh}). Listing pokazuje adres, słowo, źródło i disasemblację.

\section{Ślad wykonania}

Emulator z opcją \code{-t plik} zapisuje po każdej instrukcji jedną linię:
\begin{txt}
00000010 00a50533 x10=0000002a      pc, słowo instrukcji, zapisany rejestr = wartość
00000014 fe051ce3 -                 instrukcja nie zapisała rejestru
\end{txt}
Procesor w RTL wypisuje dokładnie ten sam format. Porównanie obu śladów
(\code{tracecmp.py}) wskazuje \emph{pierwszą} instrukcję, w której sprzęt
odbiega od modelu. To najskuteczniejsza technika debugowania procesorów.

\section{Emulator: struktura}

Gotowe są pamięć z MMIO, ładowanie \code{.hex}, pętla główna i ślad. Ty piszesz
funkcję \code{step()}:

\begin{cc}
// Wykonuje instrukcję spod c->pc. Zwraca numer zapisanego rejestru
// (0, jeśli instrukcja nic nie zapisała), wartość przez *rd_val.
static int step(cpu_t *c, uint32_t instr, uint32_t *rd_val) {
    uint32_t op  = bits(instr, 6, 0);
    uint32_t rd  = bits(instr, 11, 7);
    ...
    switch (op) {
    case 0x37: res = instr & 0xFFFFF000u; break;    // LUI (gotowe)
    // TODO: AUIPC, JAL, JALR, BRANCH, LOAD, STORE, OP-IMM, OP, FENCE, SYSTEM
    }
    if (writes && rd != 0) c->x[rd] = res;
    c->pc = next;
    ...
}
\end{cc}

Wskazówki dla C: porównanie ze znakiem to \code{(int32\_t)a < (int32\_t)b},
przesunięcie arytmetyczne to \code{(uint32\_t)((int32\_t)a >> n)} (formalnie
\emph{implementation-defined}, w praktyce działa wszędzie), a przesunięcie o
$\geq 32$ to w C UB, więc maskuj: \code{b \& 31}.

\section{Zadania}

\begin{zadanie}{4.1 --- Kodowanie na kartce}
Zakoduj ręcznie do postaci szesnastkowej i sprawdź \code{rvtool dis}:
(1) \code{addi a0, a0, -1}, (2) \code{sw ra, 12(sp)}, (3) \code{srai a1, a2, 3},
(4) \code{lui t0, 0x10000}, (5) \code{beq a0, zero, +16}, (6) \code{jal ra, -2048}.
Przy 5 i 6 rozpisz bity immediate'a na pozycje w instrukcji (rys.~\ref{fig:formaty}).
\end{zadanie}

\begin{zadanie}{4.2 --- Generator immediate'ów (\code{rtl/imm\_gen.sv})}
\code{make unit}. Wektory (400 instrukcji wszystkich formatów) wygenerował
asembler. Gdy test zgłosi błąd, zdekoduj instrukcję: \code{rvtool.py dis 0x...}.
\end{zadanie}

\begin{zadanie}{4.3 --- Emulator (\code{emu/rv32sim.c})}
\code{make emu-test} uruchamia 12 programów, a \code{make emu-test P=07\_load\_store}
jeden program ze śladem. Programy są ułożone od najprostszych (\code{01\_smoke}
potrzebuje tylko \code{lui}, \code{addi}, \code{sw}, \code{jal}), więc
implementuj instrukcje w kolejności testów. Pisz z myślą o porównywaniu z RTL:
ten emulator jest specyfikacją, więc jasność jest ważniejsza niż wydajność.
\end{zadanie}

\begin{zadanie}{4.4 --- Czytanie programu}
\code{make emu-test P=10\_recursion}, potem otwórz \plik{build/hex/10\_recursion.lst} i
\plik{build/10\_recursion.trace}. Prześledź, jak \code{sp} maleje przy wywołaniach
\code{fib}. Ile instrukcji wykonuje \code{mul(1234, 5678)} i dlaczego tyle? Ile
razy wykonało się \code{jalr} i skąd ta liczba?
\end{zadanie}

\begin{gotowe}
\code{make test} → PASS dla \code{imm\_gen} i 12 × PASS emulatora.
\end{gotowe}

\begin{kontrolne}
  \item Który bit instrukcji jest zawsze bitem znaku immediate'a i dlaczego to ułatwia sprzęt?
  \item Jaki jest zasięg \code{beq} i \code{jal}? Dlaczego bit 0 offsetu nie jest kodowany?
  \item Rozwiń \code{li a0, 0xDEADBEEF} na \code{lui} i \code{addi}.
  \item Co znajdzie się w rejestrze po \code{lb} z bajtu \code{0x80}, a co po \code{lbu}?
  \item Dlaczego emulator jest dobrym „wzorcem” dla RTL, choć sam może mieć błędy?
\end{kontrolne}

\chapter{Procesor jednocyklowy RV32I}
\label{ch:single}

\metryczka{złożyć działający procesor, który wykonuje prawdziwe programy w kodzie
maszynowym RISC-V, po jednej instrukcji na takt zegara.}{3--4 wieczory}{05-single-cycle/}

\section{Ścieżka danych}

\begin{figure}[h]
\centering
\resizebox{\linewidth}{!}{% Ścieżka danych procesora jednocyklowego (rozdział 5)
\begin{tikzpicture}[x=1cm, y=1cm, every node/.append style={font=\sffamily\scriptsize},
  szyna/.style={thick}, ster/.style={densely dashed, pulapka, thin, ->}]
  % --- bloki
  \node[draw, thick, fill=zadaniejasny, minimum width=6mm, minimum height=16mm] (pc) at (0,0) {PC};
  \node[blok, minimum width=15mm, minimum height=16mm] (imem) at (1.9,0) {pamięć\\instrukcji};
  \node[blok, minimum width=15mm, minimum height=9mm] (ctrl) at (5.4,2.75) {control};
  \node[blok, minimum width=16mm, minimum height=20mm] (rf) at (5.4,0) {regfile};
  \node[blok, minimum width=16mm, minimum height=8mm] (imm) at (5.4,-2.1) {imm\_gen};
  \node[blok, minimum width=7mm, minimum height=10mm] (ma) at (7.9,0.55) {A};
  \node[blok, minimum width=7mm, minimum height=10mm] (mb) at (7.9,-0.55) {B};
  \node[blok, minimum width=12mm, minimum height=24mm, fill=pytaniejasny] (alu) at (9.5,0) {ALU};
  \node[blok, minimum width=17mm, minimum height=8mm, fill=pytaniejasny] (br) at (9.5,-2.7) {branch\_cmp};
  \node[blok, minimum width=17mm, minimum height=16mm] (mem) at (11.9,0) {LSU +\\pamięć\\danych};
  \node[blok, minimum width=7mm, minimum height=14mm] (wb) at (13.9,0) {WB};
  \node[blok, minimum width=15mm, minimum height=9mm] (npc) at (1.9,-2.6) {następny\\PC};
  % --- PC -> imem -> instr
  \draw[szyna, ->] (pc) -- (imem);
  \draw[szyna] (imem.east) -- (3.4,0);
  \draw[szyna] (3.4,2.75) -- (3.4,-2.1);
  \node[sygnal, rotate=90, fill=white, inner sep=1pt] at (3.4,1.4) {instr};
  \draw[szyna, ->] (3.4,2.75) -- (ctrl.west);
  \draw[szyna, ->] (3.4,0.45) -- node[above, sygnal]{rs1, rs2, rd} ($(rf.west)+(0,0.45)$);
  \draw[szyna, ->] (3.4,-2.1) -- (imm.west);
  % --- regfile -> A/B -> ALU
  \draw[szyna, ->] ($(rf.east)+(0,0.55)$) -- node[above, sygnal, pos=0.75]{rs1} (ma.west);
  \draw[szyna, ->] ($(rf.east)+(0,-0.55)$) -- node[below, sygnal, pos=0.75]{rs2} (mb.west);
  \draw[szyna, ->] (imm.east) -- (7.9,-2.1) node[right, sygnal, pos=0.4, yshift=5pt]{} -- (mb.south);
  \node[sygnal] at (8.15,-1.6) {imm};
  \draw[szyna, <-] (ma.north) -- ++(0,0.45) node[left, sygnal]{pc, 0};
  \draw[szyna, ->] (ma.east) -- ($(alu.west)+(0,0.55)$);
  \draw[szyna, ->] (mb.east) -- ($(alu.west)+(0,-0.55)$);
  % --- komparator (odczepy rs1, rs2)
  \fill (6.55,0.55) circle (1.3pt);
  \fill (6.85,-0.55) circle (1.3pt);
  \draw[szyna, ->] (6.55,0.55) -- (6.55,-2.55) -- ($(br.west)+(0,0.15)$);
  \draw[szyna, ->] (6.85,-0.55) -- (6.85,-2.85) -- ($(br.west)+(0,-0.15)$);
  % --- ALU -> pamięć -> WB
  \draw[szyna, ->] (alu.east) -- node[above, sygnal]{alu\_y} (mem.west);
  \draw[szyna, ->] (mem.east) -- node[above, sygnal, pos=0.45]{load} (wb.west);
  \fill (10.55,0) circle (1.3pt);
  \draw[szyna, ->] (10.55,0) -- (10.55,-1.3) -| ($(wb.south)+(-0.15,0)$);
  \draw[szyna, <-] ($(wb.south)+(0.15,0)$) -- ++(0,-0.8) node[below, sygnal]{pc+4};
  % --- zapis z WB do regfile
  \draw[szyna, ->] (wb.east) -- ++(0.35,0) |- (5.9,2.2) -- ($(rf.north)+(0.5,0)$);
  \node[sygnal, fill=white, inner sep=1pt] at (11.2,2.2) {wynik do rejestru rd};
  % --- następny PC: imm, taken, alu_y
  \draw[szyna, ->] (imm.south) -- (5.4,-3.55) -| ($(npc.south)+(-0.35,0)$);
  \node[sygnal, fill=white, inner sep=1pt] at (4.0,-3.55) {imm};
  \draw[szyna, ->] (br.south) -- (9.5,-3.85) -| (npc.south);
  \node[sygnal, fill=white, inner sep=1pt] at (7.6,-3.85) {taken};
  \fill (10.55,-1.3) circle (1.3pt);
  \draw[szyna, ->] (10.55,-1.3) -- (10.55,-4.15) -| ($(npc.south)+(0.35,0)$);
  \node[sygnal, fill=white, inner sep=1pt] at (8.5,-4.15) {alu\_y (jalr)};
  \draw[szyna, ->] (npc.west) -- (-0.75,-2.6) |- (pc.west);
  \node[sygnal, rotate=90] at (-1.0,-1.3) {pc+4 / pc+imm / alu\_y};
  % --- sterowanie (przerywane)
  \draw[ster] (ctrl.east) -| ($(ma.north)+(0.3,0)$);
  \draw[ster] ($(ctrl.east)+(0,0.15)$) -| (alu.north);
  \draw[ster] ($(ctrl.east)+(0,0.3)$) -| (wb.north);
  \node[sygnal, text=pulapka, anchor=west] at (6.3,3.25) {sygnały sterujące (wb\_sel, alu\_op, \dots)};
\end{tikzpicture}
}
\caption{Ścieżka danych procesora jednocyklowego. Linie przerywane to sygnały
sterujące z dekodera.}
\label{fig:datapath}
\end{figure}

W jednym takcie, od zbocza do zbocza (rys.~\ref{fig:datapath}):
\begin{enumerate}
  \item \code{pc} adresuje pamięć instrukcji, a \code{instr} pojawia się kombinacyjnie.
  \item Dekoder (\code{control}) i \code{imm\_gen} analizują \code{instr}, a bank
    rejestrów czyta \code{rs1} i \code{rs2}.
  \item ALU liczy wynik, adres dla load/store albo cel skoku dla \code{jalr}.
  \item Load: LSU wybiera bajty ze słowa z pamięci danych. Store: LSU ustawia
    \code{wstrb}/\code{wdata}, a zapis nastąpi na zboczu.
  \item Multiplekser WB wybiera, co trafi do \code{rd}: wynik ALU, daną z
    pamięci albo \code{pc+4}. Zapis nastąpi na zboczu.
  \item Logika następnego PC wybiera \code{pc+4}, \code{pc+imm} albo
    \code{alu\_y \& \textasciitilde1}. Nowa wartość też trafia do PC na zboczu.
\end{enumerate}

\textbf{Wszystkie efekty instrukcji (zapis rejestru, zapis pamięci, nowe PC)
dzieją się na tym samym zboczu zegara.} Dlatego ten procesor jest taki prosty.
Cena: okres zegara musi pomieścić całą drogę od PC przez pamięć instrukcji,
rejestry, ALU i pamięć danych aż do zapisu.

\section{Przejście instrukcji przez ścieżkę danych}

\begin{center}
\small
\begin{tabular}{@{}llllll@{}}
\toprule
\textbf{Instrukcja} & \textbf{A} & \textbf{B} & \textbf{ALU} & \textbf{Do rd} & \textbf{Następne PC} \\
\midrule
\code{add rd, rs1, rs2} & rs1 & rs2 & \{i[30], f3\} & alu\_y & pc+4 \\
\code{addi rd, rs1, imm} & rs1 & imm & \{0, f3\}$^*$ & alu\_y & pc+4 \\
\code{lw rd, imm(rs1)} & rs1 & imm & ADD & load\_data & pc+4 \\
\code{sw rs2, imm(rs1)} & rs1 & imm & ADD & --- (zapis rs2) & pc+4 \\
\code{beq rs1, rs2, off} & --- & --- & --- & --- & taken ? pc+imm : pc+4 \\
\code{jal rd, off} & --- & --- & --- & pc+4 & pc+imm \\
\code{jalr rd, imm(rs1)} & rs1 & imm & ADD & pc+4 & alu\_y \& \textasciitilde1 \\
\code{lui rd, imm} & 0 & imm & ADD & alu\_y & pc+4 \\
\code{auipc rd, imm} & pc & imm & ADD & alu\_y & pc+4 \\
\bottomrule
\end{tabular}

\smallskip
{\footnotesize $^*$ dla \code{srai}: \{i[30], f3\}}
\end{center}

\section{Tabela sterowania}

Wypełnij tę tabelę \textbf{na kartce}, zanim napiszesz \code{control.sv}.
Kreska oznacza „obojętne”. Test \code{tb\_control} sprawdza dokładnie tę tabelę na 97
zakodowanych instrukcjach i wypisuje różnice z nazwami pól.

\begin{center}
\small
\renewcommand{\arraystretch}{1.6}
\begin{tabular}{|l|p{9mm}|p{9mm}|p{9mm}|p{9mm}|p{9mm}|p{9mm}|p{9mm}|p{9mm}|p{9mm}|p{9mm}|}
\hline
& \rotatebox{90}{\code{reg\_write}\ } & \rotatebox{90}{\code{wb\_sel}\ } & \rotatebox{90}{\code{alu\_a\_sel}\ } & \rotatebox{90}{\code{alu\_b\_imm}\ } & \rotatebox{90}{\code{alu\_op}\ } & \rotatebox{90}{\code{mem\_read}\ } & \rotatebox{90}{\code{mem\_write}\ } & \rotatebox{90}{\code{branch}\ } & \rotatebox{90}{\code{jump}\ } & \rotatebox{90}{\code{jump\_reg}\ } \\
\hline
\textsf{OP (R)} & & & & & & & & & & \\
\hline
\textsf{OP-IMM} & & & & & & & & & & \\
\hline
\textsf{LOAD} & & & & & & & & & & \\
\hline
\textsf{STORE} & & & & & & & & & & \\
\hline
\textsf{BRANCH} & & & & & & & & & & \\
\hline
\textsf{JAL} & & & & & & & & & & \\
\hline
\textsf{JALR} & & & & & & & & & & \\
\hline
\textsf{LUI} & & & & & & & & & & \\
\hline
\textsf{AUIPC} & & & & & & & & & & \\
\hline
\textsf{FENCE/SYSTEM} & & & & & & & & & & \\
\hline
\end{tabular}
\end{center}

\begin{pulapka}[\texttt{addi} z ujemną stałą to nie \texttt{sub}]
Dla OP-IMM \code{instr[30]} jest częścią immediate'a. Gdyby dekoder brał go do
\code{alu\_op} jak dla R-type, \code{addi a0, a0, -1} (stała z ustawionym bitem
30) liczyłoby się jako odejmowanie. Bit 30 ma znaczenie tylko dla \code{srai}
(\code{funct3 = 101}).
\end{pulapka}

\section{LSU: bajty w słowie}

Pamięć jest zorganizowana w 32-bitowe słowa. \code{addr[31:2]} wybiera słowo, a
\code{addr[1:0]} bajt w słowie. Zapis częściowy realizuje maska
\code{wstrb} (\emph{write strobe}): bit $i$ = 1 oznacza zapis bajtu $i$.

\begin{figure}[h]
\centering
\begin{tikzpicture}[x=15mm]
  \foreach \i/\r in {0/{[7:0]},1/{[15:8]},2/{[23:16]},3/{[31:24]}} {
    \draw[thick, fill=akcentjasny] (3-\i,0) rectangle ++(1,0.7);
    \node[sygnal] at (3.5-\i,0.35) {bajt \i};
    \node[sygnal, text=szary] at (3.5-\i,-0.25) {\r};
  }
  \draw[thick, fill=pulapkajasny] (1,0) rectangle ++(1,0.7);
  \node[sygnal] at (1.5,0.35) {bajt 2};
  \node[anchor=west, font=\sffamily\small] at (4.4,0.35) {\code{sb x5, 2(x0)}: \code{wstrb = 4'b0100}, \code{wdata[23:16] = x5[7:0]}};
\end{tikzpicture}
\caption{Zapis bajtu pod adres o \code{addr[1:0] = 2}.}
\end{figure}

Przy odczycie pamięć zwraca całe słowo, a LSU przesuwa je o
\code{8*addr[1:0]} i rozszerza znakiem (\code{lb}, \code{lh}) albo zerami
(\code{lbu}, \code{lhu}).

\section{Interfejs rdzenia: kontrakt z testbenchem}

\begin{sv}
module core import rv32i_pkg::*; (
  input  logic        clk, rst,
  output logic [31:0] imem_addr,  input logic [31:0] imem_rdata,
  output logic [31:0] dmem_addr,  output logic dmem_re,
  input  logic [31:0] dmem_rdata,
  output logic [3:0]  dmem_wstrb, output logic [31:0] dmem_wdata,
  output logic        commit_valid,                    // ślad wykonania
  output logic [31:0] commit_pc, commit_instr, commit_rd_val,
  output logic [4:0]  commit_rd
);
\end{sv}

\begin{itemize}
  \item Pamięć (64 KiB) jest \emph{jedna}, widziana przez dwa porty. Odczyt jest
    kombinacyjny, a zapis następuje na zboczu, gdy \code{dmem\_wstrb != 0}.
  \item \code{dmem\_addr} to pełny adres bajtowy, a \code{dmem\_rdata} całe
    słowo. Wybór bajtów robi Twoje LSU.
  \item Reset synchroniczny, a po nim \code{pc = 0}.
  \item \textbf{Interfejs commit} służy tylko do weryfikacji. W każdym takcie,
    w którym kończy się instrukcja, \code{commit\_valid = 1} i opis tej
    instrukcji. \code{commit\_rd = 0}, gdy instrukcja nie zapisuje rejestru.
\end{itemize}

Testbench (\plik{common/tb/soc\_tb.sv}) przerywa symulację z czytelnym
komunikatem, gdy: PC jest X po resecie, \code{commit\_valid} jest X, następuje
zapis pod adres X albo nieobsługiwany adres, wykona się słowo
\code{0x00000000} (skok w pustą pamięć), przez 1000 taktów nie ma commitu albo
przekroczony zostanie limit cykli.

\section{Jak debugować procesor}

\begin{enumerate}
  \item Zacznij od \textbf{najprostszego testu}, który nie przechodzi:
    \code{make prog P=01\_smoke}.
  \item \textbf{Ślad.} \code{make prog P=...} zapisuje \plik{build/P.trace}.
    Porównaj go z emulatorem:
\begin{sh}
make -C ../04-isa-emulator emu-test P=05_branch
python3 ../common/tools/tracecmp.py ../04-isa-emulator/build/05_branch.trace \
                                    build/05_branch.trace
\end{sh}
    Wynik pokazuje pierwszą różniącą się instrukcję z disasemblacją i
    podpowiedzią: inne PC oznacza błąd skoku, inna wartość błąd wykonania.
  \item \textbf{Listing} \plik{build/hex/P.lst} łączy adresy z liniami źródła.
  \item \textbf{Przebiegi}: \code{make wave P=05\_branch}. Dodaj \code{dut.pc},
    \code{imem\_rdata}, sygnały sterujące i \code{commit\_*}.
  \item „FAIL: podtest N” oznacza, że w źródle testu przed nieudanym sprawdzeniem
    stoi \code{li gp, N}.
\end{enumerate}

Przykładowy raport \code{tracecmp} dla błędu w dekoderze:

\begin{txt}
RÓŻNICA w instrukcji nr 12:
  ostatnie zgodne:
    00000024: 00500593  addi a1, zero, 5                 x11=00000005
  oczekiwano (ref): 00000028: fff58593  addi a1, a1, -1   x11=00000004
  jest       (RTL): 00000028: fff58593  addi a1, a1, -1   x11=00000006
  -> ta sama instrukcja, inny efekt: błąd wykonania (ALU, LSU, forwarding...)
\end{txt}

\section{Programy testowe}

\begin{center}
\small
\begin{tabularx}{\linewidth}{@{}l X@{}}
\toprule
\code{01\_smoke} & tylko \code{lui}, \code{addi}, \code{sw}: najprostszy możliwy test \\
\code{02\_alu\_imm}, \code{03\_alu\_reg} & ALU z immediate i rejestr-rejestr, przypadki brzegowe \\
\code{04\_lui\_auipc} & stałe górne i adresowanie względem PC \\
\code{05\_branch} & wszystkie skoki warunkowe, także dalekie ($>$2 KiB) \\
\code{06\_jal\_jalr} & adresy powrotu, \code{rd == rs1}, stos, zagnieżdżone wywołania \\
\code{07\_load\_store} & wszystkie szerokości, rozszerzanie znakiem i zerami \\
\code{08\_fib}, \code{09\_sort}, \code{10\_recursion} & „prawdziwe” programy: pętle, tablice, rekurencja \\
\code{11\_hazards} & sekwencje groźne dla potoku (rozdział 8) \\
\code{12\_mmio} & wypisuje napis przez PUTCHAR i czyta licznik cykli \\
\bottomrule
\end{tabularx}
\end{center}

Gdy przejdzie \code{12\_mmio}, zobaczysz w terminalu tekst wypisany przez
Twój procesor.

\section{Zadania}

\begin{zadanie}{5.1 --- Dekoder (\code{rtl/control.sv})}
\code{make unit T=control}.
\end{zadanie}

\begin{zadanie}{5.2 --- LSU (\code{rtl/lsu.sv})}
\code{make unit T=lsu}. Rozrysuj sobie słowo z czterema bajtami i zobacz, gdzie
trafia \code{sb} pod \code{addr[1:0] = 3}.
\end{zadanie}

\begin{zadanie}{5.3 --- Rdzeń (\code{rtl/core.sv})}
Złóż ścieżkę danych z gotowych bloków. Moduły z rozdziałów 3 i 4 są dołączane z
tamtych katalogów, więc poprawka w \plik{alu.sv} od razu działa także tutaj.
\code{make progs} uruchamia 12 programów, \code{make prog P=08\_fib} jeden, a
\code{make test} wszystko.
\end{zadanie}

\begin{zadanie}{5.4 --- Pomiary}
\code{make synth T=core}: ile komórek ma Twój procesor w porównaniu z sumą ALU
i banku rejestrów? CPI wynosi z definicji 1,000. W rozdziale 8 porównasz to z
potokiem.
\end{zadanie}

\begin{uwaga}[Uwaga o długości ścieżki]
Pamięci są w testbenchu, poza modułem \code{core}, więc Yosys widzi osobne
kawałki: PC → \code{imem\_addr}, \code{imem\_rdata} → … → \code{dmem\_addr} oraz
\code{dmem\_rdata} → rejestry. Prawdziwy takt procesora jednocyklowego to
\emph{suma}: odczyt imem + dekodowanie + rejestry + ALU + odczyt dmem + zapis.
To jest motywacja potoku.
\end{uwaga}

\begin{gotowe}
\code{make test} → 2 × PASS (testy jednostkowe) + 12 × PASS (programy), \code{make lint} czysto.
\end{gotowe}

\begin{kontrolne}
  \item Wymień wszystkie elementy, przez które przechodzi \code{lw} w ciągu jednego taktu.
  \item Po co multiplekser A ma wejścia \code{pc} i \code{0}? Które instrukcje ich używają?
  \item Dlaczego cel \code{jalr} liczy ALU, a cel \code{jal} i \code{beq} osobny sumator?
  \item Jakie \code{wstrb} i \code{wdata} wystawi LSU dla \code{sh x7, 2(x0)}, gdy \code{x7 = 0x1234BEEF}?
  \item Dlaczego procesor jednocyklowy ma niską maksymalną częstotliwość zegara?
\end{kontrolne}

\chapter{Weryfikacja: co-symulacja, fuzzing, Verilator}
\label{ch:wer}

\metryczka{przestać ufać zdaniu „przeszło 12 testów”. W prawdziwych projektach
weryfikacja zajmuje więcej czasu niż projektowanie. Ten rozdział nie ma nowego
RTL: tymi samymi narzędziami sprawdzasz rdzeń z rozdziału 5, a później potok z
rozdziału 8.}{1--2 wieczory}{06-weryfikacja/}

\section{Trzy poziomy testowania}

\begin{description}
  \item[Testy kierowane] (\emph{directed}): ręcznie napisane programy
    sprawdzające konkretne zachowania (\plik{common/sw/tests/*.S}). Dobre na
    start i na znane przypadki brzegowe, słabe na te, o których nie
    pomyślałeś.
  \item[Co-symulacja z modelem referencyjnym:] ten sam program wykonują
    emulator i RTL, a porównujemy \emph{każdą} wykonaną instrukcję (PC, słowo,
    zapisany rejestr, wartość). Błąd zostaje wykryty tam, gdzie powstał, a nie
    500 instrukcji później, gdy program „coś” źle policzy.
  \item[Testy losowe] (\emph{fuzzing}, w branży \emph{constrained random}):
    generator tworzy tysiące programów, których nikt by ręcznie nie napisał, a
    co-symulacja ocenia, czy RTL zgadza się z modelem. Szczególnie skuteczne
    przy potoku: znajdują kombinacje zależności, skoków i loadów, o których nie
    wiedziałeś, że mogą być groźne.
\end{description}

\begin{figure}[h]
\centering
\begin{tikzpicture}[node distance=7mm and 12mm, every node/.style={font=\sffamily\small}]
  \node[blok, fill=white] (gen) {\code{rvgen.py}\\(losowy)};
  \node[blok, fill=white, below=of gen] (src) {\code{tests/*.S}\\(kierowany)};
  \node[blok, right=of $(gen)!0.5!(src)$] (asm) {\code{rvtool asm}};
  \node[blok, right=of asm, yshift=9mm, fill=zadaniejasny] (emu) {emulator\\(rozdz. 4)};
  \node[blok, right=of asm, yshift=-9mm, fill=pytaniejasny] (rtl) {Twój RTL\\+ \code{soc\_tb}};
  \node[blok, right=of $(emu)!0.5!(rtl)$, fill=pulapkajasny] (cmp) {\code{tracecmp.py}};
  \node[right=8mm of cmp, align=left, font=\sffamily\scriptsize] (res) {ZGODNE\\albo pierwsza\\różnica};
  \draw[->, thick] (gen) -- (asm);
  \draw[->, thick] (src) -- (asm);
  \draw[->, thick] (asm) -- node[above left, sygnal]{.hex} (emu);
  \draw[->, thick] (asm) -- (rtl);
  \draw[->, thick] (emu) -- node[above right, sygnal]{ref.trace} (cmp);
  \draw[->, thick] (rtl) -- node[below right, sygnal]{dut.trace} (cmp);
  \draw[->, thick] (cmp) -- (res);
\end{tikzpicture}
\caption{Co-symulacja: ten sam program w dwóch niezależnych implementacjach.}
\label{fig:cosim}
\end{figure}

\paragraph{Skąd model wie, co jest poprawne?} Nie wie. Emulator też może mieć
błąd. Siła metody polega na tym, że dwie \emph{niezależne} implementacje (C i
RTL, napisane różnymi technikami) rzadko mylą się tak samo. Gdy się różnią,
jedna z nich ma błąd i rozstrzyga specyfikacja. Tak samo weryfikuje się
prawdziwe procesory: otwarte rdzenie RISC-V porównują ślad z modelami Spike albo
Sail, a programy generuje m.in. \code{riscv-dv}.

\section{Ograniczenia porównania śladów}

\begin{itemize}
  \item Ślad zawiera tylko zapisy do rejestrów. Błędny \textbf{store} wyjdzie
    dopiero przy późniejszym odczycie tego adresu. (Możliwe rozszerzenie: dodać
    do interfejsu commit adres i daną zapisu.)
  \item Odczyt licznika cykli (MMIO \code{CYCLES}) z natury daje różne wartości w
    emulatorze i w RTL. Dlatego \code{make cosim} pomija test \code{12\_mmio}, a
    ślady programów w C wywołujących \code{cycles()} rozjadą się w tym miejscu.
  \item Ślad RTL może być dłuższy od wzorcowego, bo testbench pozwala potokowi
    dokończyć instrukcje po zapisie do TOHOST. \code{tracecmp} porównuje
    długość śladu wzorcowego.
\end{itemize}

\section{Verilator}

Icarus to klasyczny symulator zdarzeniowy, który interpretuje opis i obsługuje
cztery wartości logiczne. Verilator \textbf{kompiluje} RTL do C++, symuluje dwie
wartości, myśli w cyklach zegara i jest zwykle 10--100× szybszy. Z tym samym
testbenchem (opcja \code{--timing} obsługuje \code{\#5} i \code{@(posedge)})
program \code{10\_recursion} na rozwiązaniu wzorcowym wykonywał się w Icarusie
ok.\ 0,55\,s, a w Verilatorze ok.\ 0,01\,s.

\begin{pulapka}[Verilator nie widzi X]
Niezainicjowany sygnał to dla Verilatora po prostu 0, więc błędy inicjalizacji
są ukryte. Dobra praktyka: debuguj w Icarusie, testuj masowo w Verilatorze.
\end{pulapka}

\begin{sh}
make progs SIM=verilator       # pierwsza kompilacja trwa kilka sekund
make fuzz N=500 SIM=verilator
\end{sh}

\section{Mutacje: kto pilnuje strażnika?}

Skąd wiadomo, że testy są dobre? Wstaw do rdzenia \emph{celowy} błąd i sprawdź,
czy go wykryją. To się nazywa \textbf{testowanie mutacyjne}. Przy tworzeniu kursu
w ten sposób odkryłem, że pierwsza wersja testów nie wykrywała błędu w bicie 11
immediate'a skoku warunkowego: żaden skok nie był dłuższy niż 2 KiB. Stąd
„dalekie skoki” w \code{05\_branch.S}.

\begin{center}
\small
\begin{tabularx}{\linewidth}{@{}r X l l l@{}}
\toprule
\# & \textbf{Mutant} & \textbf{progs} & \textbf{cosim} & \textbf{fuzz} \\
\midrule
1 & \code{sra} działa jak \code{srl} & & & \\
2 & \code{lh} nie rozszerza znakiem & & & \\
3 & zapis do \code{x0} nie jest ignorowany & & & \\
4 & \code{jalr} nie zeruje bitu 0 & & & \\
5 & \code{sb} zawsze zapisuje bajt 0 słowa & & & \\
6 & bit 11 B-immediate'a z \code{instr[31]} zamiast \code{instr[7]} & & & \\
7 & \code{sltiu} porównuje ze znakiem & & & \\
\bottomrule
\end{tabularx}
\end{center}

\section{Pokrycie funkcjonalne}

„Ile testów przeszło” to zła miara. Lepiej pytać, \emph{które zachowania zostały
sprawdzone}. Czy każda instrukcja się wykonała? Czy \code{bgeu} był choć raz
\emph{wzięty}, a nie tylko niewzięty? Czy wystąpił load tuż po store pod ten sam
adres? Tak myśli się o pokryciu (\emph{coverage}). Narzędzia przemysłowe zbierają
je automatycznie (\code{covergroup} w SystemVerilogu). W kursie policzysz je
skryptem ze śladów.

\section{Zadania}

Polecenia uruchamiasz z katalogu \plik{06-weryfikacja}. Domyślnie testowany jest
rdzeń z rozdziału 5, a z \code{CORE=08} potok z rozdziału 8.

\begin{zadanie}{6.1 --- Co-symulacja}
\code{make cosim} wymaga działającego emulatora i rdzenia. Jeśli coś się
rozjedzie, przeczytaj raport \code{tracecmp} i znajdź błąd. Pamiętaj, że może
być także w emulatorze!
\end{zadanie}

\begin{zadanie}{6.2 --- Fuzzing}
\code{make fuzz} (20 programów po 300 instrukcji), potem
\code{make fuzz N=300 LEN=600 SEED=1000}. Przeczytaj \plik{common/tools/rvgen.py}
i jeden wygenerowany program (\plik{build/core05/fuzz/r1.S}). Czego generator
\emph{nie} tworzy?
\end{zadanie}

\begin{zadanie}{6.3 --- Polowanie na mutanty}
Wprowadzaj do rdzenia po jednym błędzie z tabeli powyżej i wypełnij ją: który
mechanizm wykrył błąd? Dla mutanta, którego fuzzer nie wykrywa, rozszerz
\code{rvgen.py} tak, żeby wykrywał (np. dalekie skoki: \code{.space} między
skokiem a celem). Na koniec cofnij mutację!
\end{zadanie}

\begin{zadanie}{6.4 --- Pokrycie}
Napisz krótki skrypt w Pythonie, który z plików \plik{build/core05/fuzz/*.trace}
policzy, ile razy wykonała się każda instrukcja (disasembler:
\code{from rvtool import dis}). Czy wystąpiło wszystkie 37 instrukcji RV32I (bez
\code{fence/ecall/ebreak})?
\end{zadanie}

\begin{zadanie}{6.5 --- Lint i synteza całego rdzenia}
\code{make lint}, \code{make synth T=core}.
\end{zadanie}

\begin{gotowe}
\code{make cosim} → wszystko ZGODNE, \code{make fuzz N=200} → 0 różnic, a
fuzzer po Twojej zmianie wykrywa mutanta nr 6.
\end{gotowe}

\begin{kontrolne}
  \item Dlaczego co-symulacja lokalizuje błąd lepiej niż test, który na końcu sprawdza wynik?
  \item Jakie błędy przejdą niezauważone przez porównanie śladu złożonego tylko z zapisów do rejestrów?
  \item Kiedy wolisz Icarusa, a kiedy Verilatora?
  \item Co mierzy testowanie mutacyjne i czym różni się od pokrycia kodu?
\end{kontrolne}

\chapter{Programy w C na Twoim procesorze}
\label{ch:c}

\metryczka{uruchomić skompilowany kod C na procesorze, który sam zbudowałeś, bez
systemu operacyjnego i bez libc. To terytorium, które znasz z pracy, tylko
tym razem pod spodem jest Twój własny sprzęt.}{1--2 wieczory}{07-c-bare-metal/}

\section{Toolchain}

Makefile wykrywa toolchain sam. Wystarczy jedno z dwóch:
\begin{sh}
sudo apt install lld                      # clang już jest, brakuje tylko linkera
sudo apt install gcc-riscv64-unknown-elf  # albo pełny toolchain GNU (umie też rv32)
\end{sh}
\code{make check-tools} pokazuje, czego brakuje. Moduł został sprawdzony z
clang~18 i \code{ld.lld}~18: wszystkie programy przechodzą na emulatorze, rdzeniu
jednocyklowym i potoku. Ścieżka GCC jest w Makefile, ale nie była testowana.

\section{Od C do pliku \texttt{.hex}}

\begin{figure}[h]
\centering
\begin{tikzpicture}[node distance=17mm, every node/.style={font=\sffamily\small}]
  \node[blok, fill=white] (c) {\code{hello.c}\\\code{rt.c}, \code{muldiv.c}};
  \node[blok, fill=white, below=4mm of c] (s) {\code{crt0.S}};
  \node[blok, right=20mm of $(c)!0.5!(s)$] (o) {\code{*.o}};
  \node[blok, right=of o] (elf) {\code{.elf}};
  \node[blok, right=of elf] (bin) {\code{.bin}};
  \node[blok, right=of bin, fill=zadaniejasny] (hex) {\code{.hex}};
  \node[blok, fill=white, above=6mm of elf] (ld) {\code{link.ld}};
  \draw[->, thick] (c) -- (o);
  \node[sygnal, above=1mm of o] {clang -c};
  \draw[->, thick] (s) -- (o);
  \draw[->, thick] (o) -- node[above, sygnal]{ld.lld} (elf);
  \draw[->, thick] (ld) -- (elf);
  \draw[->, thick] (elf) -- node[above, sygnal]{objcopy} (bin);
  \draw[->, thick] (bin) -- node[above, sygnal]{bin2hex} (hex);
\end{tikzpicture}
\caption{Kompilacja programu na nasz procesor.}
\end{figure}

Najważniejsze flagi:
\begin{itemize}
  \item \code{-march=rv32i -mabi=ilp32}: tylko instrukcje bazowe;
    \code{int}, \code{long} i wskaźnik mają po 32 bity,
  \item \code{-ffreestanding -nostdlib}: nie ma libc ani standardowego kodu startowego,
  \item \code{-mno-relax}: linker nie zamienia adresowania na względne
    wobec \code{gp} (nasz crt0 nie ustawia \code{gp}),
  \item \code{-O2}: warto zobaczyć, jak dobry kod generuje kompilator
    (\code{make dis P=bench}).
\end{itemize}

\section{Co się dzieje przed \texttt{main()}}

Procesor po resecie zaczyna od \code{pc = 0} i nic nie wie o C. Kod startowy
\plik{crt0.S} musi:
\begin{asm}
_start:
    la   sp, _stack_top        # 1. stos: C go potrzebuje do zmiennych lokalnych
    la   t0, __bss_start       # 2. zerowanie .bss: "static int x;" ma być 0
    la   t1, __bss_end
1:  bgeu t0, t1, 2f
    sw   zero, 0(t0)
    addi t0, t0, 4
    j    1b
2:  call main                  # 3. program
    ...                        # 4. a0 == 0 ? TOHOST = 1 : TOHOST = (a0 << 1) | 1
\end{asm}

Sekcji \code{.data} nie trzeba kopiować, bo cały obraz programu ląduje w RAM przez
\code{\$readmemh}. Na prawdziwym mikrokontrolerze \code{.data} leży we flashu i
crt0 kopiuje ją do RAM.

\section{Skrypt linkera}

\begin{figure}[h]
\centering
\begin{tikzpicture}[x=1cm, y=1cm]
  \foreach \y/\h/\n/\c in {0/0.9/{.text (crt0 najpierw!)}/akcentjasny, 0.9/0.6/.rodata/akcentjasny,
                         1.5/0.5/.data/zadaniejasny, 2.0/0.5/.bss/zadaniejasny,
                         2.5/1.6/{wolne}/white, 4.1/0.9/{stos $\downarrow$}/pulapkajasny} {
    \draw[thick, fill=\c] (0,-\y) rectangle (5,-\y-\h);
    \node[font=\sffamily\small] at (2.5,-\y-\h/2) {\n};
  }
  \node[sygnal, anchor=east] at (-0.1,0) {0x0000};
  \node[sygnal, anchor=east] at (-0.1,-2.0) {\_\_bss\_start};
  \node[sygnal, anchor=east] at (-0.1,-2.5) {\_\_bss\_end};
  \node[sygnal, anchor=east] at (-0.1,-5.0) {0x10000 = \_stack\_top};
\end{tikzpicture}
\caption{Układ pamięci programu (64 KiB RAM).}
\end{figure}

Symbole \code{\_\_bss\_start}, \code{\_\_bss\_end} i \code{\_stack\_top} definiuje
\plik{link.ld}, a używa ich \plik{crt0.S}. W ten sposób linker przekazuje adresy do
kodu.

\section{Brakujące instrukcje: \texttt{\_\_mulsi3} i spółka}

RV32I nie ma mnożenia ani dzielenia. Gdy piszesz \code{a * b}, kompilator wstawia
\code{call \_\_mulsi3}. Zwykle dostarcza ją \code{libgcc} albo
\code{compiler-rt}, a u nas piszesz ją sam (zadanie 7.2). Zasady:

\begin{itemize}
  \item w \plik{muldiv.c} \textbf{nie wolno} używać \code{*}, \code{/} ani
    \code{\%}, bo kompilator wywołałby te same funkcje i powstałaby nieskończona
    rekurencja,
  \item dzielenie przez zero ma dawać to samo co instrukcje RISC-V M:
    $x / 0 = $ \code{0xFFFFFFFF}, $x \bmod 0 = x$,
  \item dzielenie ze znakiem zaokrągla do zera: $-7/2 = -3$, $-7 \bmod 2 = -1$.
\end{itemize}

Mnożenie podaję jako gotowy przykład, a dzielenie piszesz sam. Mnożenie „pisemne” w systemie dwójkowym:
\begin{cc}
uint32_t __mulsi3(uint32_t a, uint32_t b) {
    uint32_t r = 0;
    while (b) {
        if (b & 1) r += a;     // bit mnożnika ustawiony: dodaj przesuniętą mnożną
        a <<= 1;
        b >>= 1;
    }
    return r;
}
\end{cc}
Dzielenie robi się analogicznie (\emph{restoring division}): przesuwaj resztę w
lewo, dosuwaj kolejny bit dzielnej, a gdy reszta $\geq$ dzielnik, odejmij i ustaw
bit ilorazu.

Kompilator potrafi też sam wstawić wywołania \code{memset} i \code{memcpy} (np.
przy zerowaniu tablicy na stosie), nawet jeśli ich nie używasz. Dlatego są w
\plik{lib/rt.c}.

\section{MMIO i \texttt{volatile}}

\begin{cc}
#define MMIO_PUTCHAR ((volatile uint32_t *)0x10000004)
*MMIO_PUTCHAR = 'A';
\end{cc}
Bez \code{volatile} kompilator mógłby uznać dwa zapisy pod ten sam adres za zbędne
i usunąć jeden. Na tym poziomie nie ma różnicy między zmienną a rejestrem
urządzenia: to ten sam \code{sw}, tylko pod inny adres, który testbench (albo
prawdziwy dekoder adresów) kieruje do urządzenia.

\section{Konwencja wywołań (ilp32)}

Argumenty trafiają do \code{a0}--\code{a7}, wynik do \code{a0}, adres powrotu do
\code{ra}. Rejestry \code{s0}--\code{s11} funkcja musi zachować, więc zapisuje je
na stosie w prologu. \code{t*} i \code{a*} może niszczyć. Stos jest wyrównany do
16 bajtów. Typowy prolog i epilog:
\begin{asm}
    addi sp, sp, -16
    sw   ra, 12(sp)
    sw   s0, 8(sp)
    ...
    lw   s0, 8(sp)
    lw   ra, 12(sp)
    addi sp, sp, 16
    ret
\end{asm}

\section{Zadania}

\begin{zadanie}{7.1 --- Hello, world}
\code{make} (kompilacja), \code{make emu P=hello} (emulator), \code{make run P=hello}
(Twój rdzeń; \code{CORE=08} uruchamia potok). Pod jakim adresem leży
\code{counter} z \plik{hello.c} i w której jest sekcji? Co by się stało bez
zerowania \code{.bss}? Ile instrukcji wykonuje \code{hello}, a ile z nich to crt0?
\end{zadanie}

\begin{zadanie}{7.2 --- \code{lib/muldiv.c}}
\code{make emu P=muldiv\_test}, potem \code{make run P=muldiv\_test}. Wektory
testowe obejmują dzielenie przez zero, \code{INT\_MIN / -1} i liczby ujemne.
\end{zadanie}

\begin{zadanie}{7.3 --- Benchmark}
\code{make run P=bench} wypisuje wyniki i liczbę cykli dla CRC32, sita
Eratostenesa, sortowania i mnożenia macierzy. Która część jest najdroższa?
Sprofiluj ją na poziomie sprzętu:
\begin{sh}
cut -d' ' -f1 build/bench.trace | sort | uniq -c | sort -rn | head
\end{sh}
i znajdź najczęstsze adresy w \plik{build/bench.dis}.
\end{zadanie}

\begin{zadanie}{7.4 --- Własny program}
Napisz \plik{prog/coś.c}: liczby pierwsze do 10\,000, drzewo BST w statycznej
tablicy albo „grę w życie” wypisywaną znakami. \code{main} zwraca 0 jako PASS.
\end{zadanie}

\begin{gotowe}
\code{make test} → wszystkie PASS na emulatorze i na rdzeniu.
\end{gotowe}

\begin{kontrolne}
  \item Co musi zrobić crt0, zanim wywoła \code{main}? Czego nie musi robić u nas, a musi na mikrokontrolerze?
  \item Skąd crt0 zna adresy początku i końca \code{.bss}?
  \item Dlaczego w \plik{muldiv.c} nie można użyć operatora \code{*}?
  \item Co się stanie, gdy usuniesz \code{volatile} z definicji \code{MMIO\_PUTCHAR} i skompilujesz z \code{-O2}?
\end{kontrolne}

\chapter{Procesor potokowy}
\label{ch:pipe}

\metryczka{przerobić procesor jednocyklowy na klasyczny 5-etapowy potok z obsługą
wszystkich hazardów i zmierzyć, co to dało. Najtrudniejszy i najciekawszy
rozdział kursu.}{4--6 wieczorów}{08-pipeline/}

% Diagram potoku: \instr{wiersz}{pierwszy takt}{etap1,etap2,...}{mnemonik}
\newcommand{\instr}[4]{%
  \node[anchor=east, font=\ttfamily\scriptsize] at (0.9,-#1) {#4};
  \foreach \s [count=\k from 0] in {#3} {
    \ifx\s\empty\else
    \pgfmathsetmacro{\xx}{#2+\k}
    \node[draw, thick, minimum width=9.5mm, minimum height=5mm, inner sep=0pt,
      font=\sffamily\scriptsize, fill=akcentjasny] at (\xx,-#1) {\s};
    \fi
  }}
\newcommand{\banka}[2]{% bańka: wiersz, takt
  \node[draw, thick, dashed, minimum width=9.5mm, minimum height=5mm, inner sep=0pt,
    font=\sffamily\tiny, fill=white, text=szary] at (#2,-#1) {bańka};}
\newcommand{\takty}[1]{%
  \foreach \t in {1,...,#1} { \node[font=\sffamily\scriptsize, text=szary] at (\t+1,0.6) {\t}; }}

\section{Po co potok?}

Czas wykonania programu („żelazne prawo” wydajności procesora):
\[ \text{czas} \;=\; \text{liczba instrukcji} \times \text{CPI} \times T_{clk}. \]

Procesor jednocyklowy ma CPI = 1, ale bardzo długi takt. Potok tnie drogę
instrukcji rejestrami na 5 krótszych etapów. Każda instrukcja nadal potrzebuje 5
etapów, więc \emph{opóźnienie} pojedynczej instrukcji się nie zmienia. Ale \textbf{w
każdym takcie kończy się jedna instrukcja}, a takt może być kilka razy krótszy.

\begin{figure}[h]
\centering
\begin{tikzpicture}[x=10.5mm, y=7mm]
  \takty{8}
  \instr{1}{2}{IF,ID,EX,MEM,WB}{add x1,..}
  \instr{2}{3}{IF,ID,EX,MEM,WB}{sub x2,..}
  \instr{3}{4}{IF,ID,EX,MEM,WB}{lw x3,..}
  \instr{4}{5}{IF,ID,EX,MEM,WB}{beq ..}
  \draw[pulapka, thick, rounded corners] (5.5,0.3) rectangle (6.5,-4.4);
  \node[etykieta, text=pulapka, align=center] at (6,-4.9) {w takcie 5 pracują\\wszystkie etapy};
\end{tikzpicture}
\caption{Idealny potok: od piątego taktu kończy się jedna instrukcja na takt.}
\label{fig:pipe}
\end{figure}

Cena: instrukcje „zachodzą” na siebie i zależności między nimi trzeba obsłużyć
w sprzęcie. To są \textbf{hazardy}.

\section{Etapy i rejestry potoku}

\begin{center}
\small
\begin{tabularx}{\linewidth}{@{}l X X@{}}
\toprule
\textbf{Etap} & \textbf{Co robi} & \textbf{Rejestr na końcu etapu} \\
\midrule
IF & \code{imem\_addr = pc}, pobranie instrukcji, \code{pc += 4} & IF/ID: \code{d\_valid, d\_pc, d\_instr} \\
ID & dekoder, immediate, odczyt rejestrów, wykrycie load-use & ID/EX: sterowanie, \code{rs1\_val, rs2\_val, imm}, numery rs1/rs2/rd \\
EX & forwarding, ALU, \code{branch\_cmp}, cel skoku, \textbf{redirect} & EX/MEM: wynik, dana do store, rd \\
MEM & LSU i pamięć danych & MEM/WB: wartość do zapisu, rd \\
WB & zapis do banku rejestrów, \textbf{commit} & --- \\
\bottomrule
\end{tabularx}
\end{center}

\begin{figure}[h]
\centering
\begin{tikzpicture}[x=1cm, y=1cm, every node/.append style={font=\sffamily\scriptsize}]
  \foreach \x/\n in {0/IF, 2.9/ID, 5.8/EX, 8.7/MEM, 11.6/WB} {
    \node[blok, minimum width=17mm, minimum height=14mm] (\n) at (\x,0) {\n};
  }
  \foreach \x/\n/\r in {1.45/{IF/ID}/ra, 4.35/{ID/EX}/rb, 7.25/{EX/MEM}/rc, 10.15/{MEM/WB}/rd} {
    \node[draw, thick, fill=zadaniejasny, minimum width=3.5mm, minimum height=22mm] (\r) at (\x,0) {};
    \node[etykieta] at (\x,1.4) {\n};
  }
  \draw[->, thick] (IF) -- (ra); \draw[->, thick] (ra) -- (ID);
  \draw[->, thick] (ID) -- (rb); \draw[->, thick] (rb) -- (EX);
  \draw[->, thick] (EX) -- (rc); \draw[->, thick] (rc) -- (MEM);
  \draw[->, thick] (MEM) -- (rd); \draw[->, thick] (rd) -- (WB);
  % redirect
  \draw[->, thick, pulapka] (EX.north) -- ++(0,1.25) -| node[above, pos=0.25]{redirect (skok wzięty)} (IF.north);
  % forwarding
  \draw[->, thick, zadanie] ($(rc.south)$) -- ++(0,-0.6) -| node[below, pos=0.3]{EX/MEM → EX} ($(EX.south)+(0.3,0)$);
  \draw[->, thick, zadanie] ($(rd.south)$) -- ++(0,-1.1) -| node[below, pos=0.2]{MEM/WB → EX} ($(EX.south)+(-0.3,0)$);
  % zapis rejestru
  \draw[->, thick, akcent] (WB.north) -- ++(0,1.85) -| node[above, pos=0.25]{zapis rd (+ bypass WB → ID)} (ID.north);
\end{tikzpicture}
\caption{Pięć etapów, cztery rejestry potoku. Czerwony: redirect przy skoku.
Zielone: ścieżki forwardingu. Niebieski: zapis do banku rejestrów.}
\label{fig:pipeline}
\end{figure}

\paragraph{Konwencje} (zgodne ze szkieletem \plik{rtl/core.sv}): przedrostek
\code{d\_} oznacza sygnał w ID, \code{e\_} w EX, \code{m\_} w MEM, a \code{w\_} w WB.
Każdy etap ma bit \code{*\_valid}. \textbf{Bańka} (\emph{bubble}) to etap z
\code{valid = 0} i wyzerowanymi sygnałami, które mają skutki uboczne
(\code{reg\_write}, \code{mem\_write}, \code{branch}, \code{jump} \dots). Przepływa
przez potok jak \code{nop}, ale nie jest raportowana jako commit.

\paragraph{Commit z WB.} \code{commit\_valid = w\_valid}, a opis instrukcji
(\code{pc}, \code{instr}) jedzie przez cały potok. Dzięki temu ślad jest identyczny
jak w emulatorze i \code{make cosim} działa bez zmian.

\section{Hazardy danych}

\begin{figure}[h]
\centering
\begin{tikzpicture}[x=10.5mm, y=7mm]
  \takty{6}
  \instr{1}{2}{IF,ID,EX,MEM,WB}{add x1,x2,x3}
  \instr{2}{3}{IF,ID,EX,MEM,WB}{sub x4,x1,x5}
  \draw[->, very thick, pulapka] (6,-1.3) .. controls (5.5,-1.7) and (4.6,-1.6) .. (4.1,-1.75);
  \node[etykieta, text=pulapka, anchor=west] at (7.7,-1) {x1 zapisany na końcu WB (takt 5)};
  \node[etykieta, text=pulapka, anchor=west] at (7.7,-2) {a czytany w ID już w takcie 3!};
\end{tikzpicture}
\caption{Hazard RAW (\emph{read after write}).}
\end{figure}

Rozwiązania, od najprostszego:
\begin{enumerate}
  \item \textbf{Programowo}: \code{nop}-y między zależnymi instrukcjami. Tak
    działa \code{make progs PAD=n} (\code{rvtool --pad-nops}) na etapie 8.1.
  \item \textbf{Bypass WB → ID}: gdy WB w tym samym takcie zapisuje rejestr, który
    ID czyta, weź wartość prosto z WB (bank rejestrów oddałby starą).
  \item \textbf{Forwarding do EX}: wynik \code{add} istnieje już na końcu EX.
    Następna instrukcja potrzebuje go na \emph{początku} swojego EX, czyli takt
    później. Przekaż go z rejestru EX/MEM albo MEM/WB prosto na wejście ALU.
  \item \textbf{Stall} przy \emph{load-use} (niżej).
\end{enumerate}

\begin{figure}[h]
\centering
\begin{tikzpicture}[x=10.5mm, y=7mm]
  \takty{6}
  \instr{1}{2}{IF,ID,EX,MEM,WB}{add x1,x2,x3}
  \instr{2}{3}{IF,ID,EX,MEM,WB}{sub x4,x1,x5}
  \instr{3}{4}{IF,ID,EX,MEM,WB}{or  x6,x1,x7}
  \draw[->, very thick, zadanie] (4.45,-1.2) -- (4.55,-1.8);
  \draw[->, very thick, zadanie] (5.45,-1.2) .. controls (5.6,-2.2) .. (5.55,-2.8);
  \node[etykieta, text=zadanie, anchor=west] at (8.8,-1.6) {EX/MEM → EX (odległość 1)};
  \node[etykieta, text=zadanie, anchor=west] at (8.8,-2.6) {MEM/WB → EX (odległość 2)};
\end{tikzpicture}
\caption{Forwarding: wynik z rejestru potoku wraca prosto na wejście ALU.}
\end{figure}

\begin{sv}
always_comb begin
  if (m_valid && m_reg_write && m_rd != 0 && m_rd == e_rs1)  fwd_a = m_result; // najnowszy!
  else if (w_valid && w_reg_write && w_rd != 0 && w_rd == e_rs1) fwd_a = w_val;
  else                                                     fwd_a = e_rs1_val;
  // to samo dla fwd_b; fwd_b jest też daną dla store!
end
\end{sv}

\begin{pulapka}[Trzy szczegóły forwardingu]
(1) \textbf{Kolejność}: gdy oba etapy piszą ten sam rejestr, wygrywa
\emph{nowszy}, czyli EX/MEM (test \code{11\_hazards}, podtest 4).
(2) \textbf{\code{rd != 0}}: zapis do \code{x0} nie może być forwardowany
(podtest 8). (3) Dana do \textbf{store} (\code{rs2}) też musi przejść przez
forwarding.
\end{pulapka}

\subsection*{Load-use: jedyny przypadek, gdy trzeba czekać}

\begin{figure}[h]
\centering
\begin{tikzpicture}[x=10.5mm, y=7mm]
  \takty{7}
  \instr{1}{2}{IF,ID,EX,MEM,WB}{lw  x1,0(x2)}
  \instr{2}{3}{IF,ID,ID,EX,MEM,WB}{add x3,x1,x1}
  \banka{3}{5}
  \node[etykieta, anchor=west] at (5.6,-3) {bańka wstawiona do ID/EX};
  \draw[->, very thick, zadanie] (5.45,-1.2) -- (5.55,-1.8);
  \node[etykieta, text=zadanie, anchor=west] at (8.1,-1.4) {dana z pamięci jest dopiero po MEM};
\end{tikzpicture}
\caption{Load-use: instrukcja zależna czeka jeden takt w ID, a potem dostaje daną przez forwarding z MEM/WB.}
\end{figure}

Wykrywasz to w ID: w EX jest load, a jego \code{rd} to \code{rs1} albo \code{rs2}
instrukcji w ID. Wtedy \textbf{PC i IF/ID stoją}, a do ID/EX idzie bańka.

\section{Hazardy sterowania}

Skok rozstrzyga się w EX. Do tego czasu IF i ID zdążyły pobrać dwie instrukcje
„za skokiem”. Jeśli skok jest wzięty: \code{pc = cel}, a instrukcje w IF/ID i
ID/EX trzeba \textbf{unieważnić} (\emph{flush}).

\begin{figure}[h]
\centering
\begin{tikzpicture}[x=10.5mm, y=7mm]
  \takty{7}
  \instr{1}{2}{IF,ID,EX,MEM,WB}{beq (wzięty)}
  \instr{2}{3}{IF,ID}{inst+4}
  \instr{3}{4}{IF}{inst+8}
  \instr{4}{5}{IF,ID,EX,MEM,WB}{cel skoku}
  \draw[pulapka, very thick] (2.55,-1.75) -- (3.45,-2.25) (2.55,-2.25) -- (3.45,-1.75);
  \draw[pulapka, very thick] (3.55,-1.75) -- (4.45,-2.25) (3.55,-2.25) -- (4.45,-1.75);
  \draw[pulapka, very thick] (3.55,-2.75) -- (4.45,-3.25) (3.55,-3.25) -- (4.45,-2.75);
  \node[etykieta, text=pulapka, anchor=west] at (6.8,-2.5) {unieważnione: kara 2 takty};
\end{tikzpicture}
\caption{Wzięty skok: dwie instrukcje pobrane „na ślepo” zamieniają się w bańki.}
\end{figure}

Strategia „zakładaj, że skok nie będzie wzięty” to najprostsza predykcja skoków.
\textbf{Priorytety}: gdy w jednym takcie jest i redirect, i stall, wygrywa redirect,
bo czekająca instrukcja i tak zostanie unieważniona.

\section{Plan pracy}

Szkielet \plik{rtl/core.sv} ma sekcje dla etapów. Pracuj etapami i commituj po
każdym, który działa.

\begin{center}
\small
\begin{tabularx}{\linewidth}{@{}l X l@{}}
\toprule
\textbf{Etap} & \textbf{Co dodajesz} & \textbf{Test} \\
\midrule
8.1 & rejestry potoku, bity \code{valid}, redirect i flush przy skokach & \code{make progs PAD=4} \\
8.2 & forwarding EX/MEM→EX i MEM/WB→EX (oba operandy i dana do store), bypass WB→ID & \code{make progs PAD=1} \\
8.3 & stall load-use & \code{make progs} \\
8.4 & weryfikacja & \code{make cosim}, \code{make fuzz N=200} \\
8.5 & pomiary: CPI, synteza & \code{make synth T=core} \\
\bottomrule
\end{tabularx}
\end{center}

\code{PAD=n} wstawia $n$ instrukcji \code{nop} \emph{przed} każdą instrukcją. Etykiety
wskazują na pierwszy z nich, żeby adresy powrotu \code{jal} się zgadzały. Przy
\code{PAD=4} żadna zależność danych nie jest groźna.

\section{Zadania}

\begin{zadanie}{8.1 --- Potok bez obsługi hazardów danych}
\code{make progs PAD=4}. Jaki jest najmniejszy PAD, przy którym Twój potok z
etapu 8.1 przechodzi wszystkie testy? Wylicz to z diagramu potoku, zanim
sprawdzisz. Jak zmienia się ta liczba po dodaniu samego bypassu WB→ID?
\end{zadanie}

\begin{zadanie}{8.2 --- Forwarding i bypass}
\code{make progs PAD=1}. Dlaczego przy \code{PAD=0} nadal coś pada? Które testy i
dlaczego akurat te?
\end{zadanie}

\begin{zadanie}{8.3 --- Stall load-use}
\code{make progs}: 12 × PASS bez żadnych \code{nop}-ów.
\end{zadanie}

\begin{zadanie}{8.4 --- Weryfikacja}
\code{make cosim}, \code{make fuzz N=200}, a także
\code{make -C ../06-weryfikacja fuzz CORE=08 N=1000 SIM=verilator} i
\code{make -C ../07-c-bare-metal test CORE=08}. Fuzzer jest tu naprawdę cenny:
losowe programy tworzą kombinacje zależności, skoków i loadów, których nie ma w
ręcznych testach.
\end{zadanie}

\begin{zadanie}{8.5 --- Pomiary}
Porównaj CPI z \code{make progs} z rozdziałem 5 (zawsze 1,000). Na rozwiązaniu
wzorcowym \code{09\_sort} ma CPI ok.\ 1,5, a \code{bench} z rozdziału 7 ok.\ 1,3. Ile
taktów zabierają skoki, a ile load-use? Potem \code{make synth T=core} tu i w
rozdziale 5. Który etap ma teraz najdłuższą ścieżkę? (Podpowiedź: EX, czyli
forwarding, ALU, decyzja o skoku i multiplekser PC.)
\end{zadanie}

\begin{zadanie}{8.6 --- Dla ambitnych}
Skoki rozstrzygane w ID (kara 1 takt, ale potrzebny forwarding do ID i nowy
stall); statyczna predykcja BTFN (skoki wstecz zakładamy jako wzięte);
predyktor dynamiczny z 2-bitowymi licznikami i BTB. Mierz CPI na \code{bench}.
\end{zadanie}

\begin{gotowe}
\code{make progs} → 12 × PASS przy \code{PAD=0}, \code{make cosim} → wszystko
ZGODNE, \code{make fuzz N=500} → 0 różnic, a \code{bench} z rozdziału 7 działa na
\code{CORE=08}.
\end{gotowe}

\begin{kontrolne}
  \item Dlaczego potok skraca czas wykonania programu, choć każda instrukcja trwa tyle samo albo dłużej?
  \item Narysuj diagram potoku dla \code{lw x1,0(x2); sw x1,4(x2)}. Czy potrzebny jest stall?
  \item Dlaczego przy forwardingu EX/MEM ma pierwszeństwo przed MEM/WB?
  \item Ile taktów traci potok na wziętym skoku i jak to zmniejszyć?
  \item Co się stanie, jeśli przy stallu zapomnisz zatrzymać rejestr IF/ID (tylko PC)?
\end{kontrolne}

\chapter{Co dalej: projekty rozszerzające}
\label{ch:dalej}

\metryczka{wybrać kierunek dalszej pracy. Masz zweryfikowany procesor w dwóch
mikroarchitekturach, emulator, infrastrukturę testową i toolchain C. Projekty
poniżej są niezależne i wszystkie korzystają z narzędzi kursu.}{bez ograniczeń}{09-dalej/}

\section{A. Rozszerzenie M (mnożenie i dzielenie)}
\emph{Około 2 wieczorów.} Instrukcje \code{mul, mulh, mulhsu, mulhu, div, divu,
rem, remu} w sprzęcie i porównanie wydajności z programowym \code{\_\_mulsi3}.
\begin{enumerate}
  \item Emulator: przypadek \code{funct7 = 0000001} dla opcode OP. Test:
    \code{make -C ../04-isa-emulator emu-test-m}. Uwaga na dzielenie przez zero i
    \code{INT\_MIN / -1}.
  \item RTL w wersji prostej: operator \code{*} w ALU (1 takt). \code{make synth}
    pokaże, ile kosztuje mnożarka.
  \item Dzielenie iteracyjne w 32 taktach: potok musi czekać (stall z EX), aż
    dzielnik skończy. To dobre ćwiczenie z wielocyklowych jednostek wykonawczych.
  \item Testy: \plik{common/sw/tests\_m/m\_ext.S}, \code{make fuzz FUZZ\_FLAGS=--mext},
    C z \code{ARCH=rv32im}, porównanie cykli \code{bench}.
\end{enumerate}

\section{B. CSR, wyjątki i przerwania}
\emph{Około 4--6 wieczorów.} Celem jest procesor, na którym da się uruchomić RTOS.
\begin{enumerate}
  \item Zicsr: \code{csrrw/csrrs/csrrc} i rejestry \code{mcycle}, \code{minstret},
    \code{mstatus}, \code{mtvec}, \code{mepc}, \code{mcause}, \code{mie}, \code{mip}.
  \item Wyjątki: \code{ecall}, \code{ebreak}, nielegalna instrukcja (sygnał
    \code{illegal} z dekodera już jest!), niewyrównany dostęp. Skok do \code{mtvec},
    powrót przez \code{mret}.
  \item W potoku wyjątek musi być \textbf{precyzyjny}: starsze instrukcje się
    kończą, a młodsze zostają unieważnione.
  \item Przerwanie od timera (\code{mtime}/\code{mtimecmp} jako MMIO).
\end{enumerate}

\section{C. Procesor na FPGA}
\emph{Około 3--5 wieczorów i płytka.} Program w C wypisuje tekst z Twojego procesora
przez UART do terminala na komputerze.

Płytki z w pełni otwartym toolchainem (Yosys + nextpnr): \textbf{iCEBreaker} lub
\textbf{iCESugar} (iCE40UP5K), \textbf{Tang Nano 9K/20K} (Gowin, Apicula),
\textbf{ULX3S} lub \textbf{OrangeCrab} (ECP5).

Co trzeba zmienić:
\begin{itemize}
  \item \textbf{Pamięć}: blokowy RAM w FPGA ma odczyt \emph{synchroniczny}. Procesor
    jednocyklowy tego wprost nie zniesie, więc potrzebny jest odczyt na zboczu
    opadającym, wersja wielocyklowa albo potok z adresem imem podawanym takt
    wcześniej.
  \item \textbf{UART}: \code{uart\_tx} z rozdziału 2 jako urządzenie MMIO z bitem
    \code{busy} do odczytu.
  \item \textbf{Zegar i reset}: PLL albo dzielnik, reset z przycisku przez
    synchronizator.
  \item \textbf{Timing}: nextpnr poda prawdziwą maksymalną częstotliwość, czyli
    rzetelne porównanie procesora jednocyklowego z potokiem.
\end{itemize}

\section{D. Pamięć z opóźnieniem i cache}
\emph{Około 4 wieczorów.} Interfejs \code{req/valid/ready} zamiast „magicznej”
pamięci, stall całego potoku przy chybieniu, cache instrukcji direct-mapped z
liniami 16 B i pomiar trafień na \code{bench}.

\section{E. Weryfikacja formalna}
\emph{Około 2--3 wieczorów.} \textbf{riscv-formal} z SymbiYosys dowodzi zgodności z
ISA. Rdzeń wystawia interfejs \textbf{RVFI}, czyli rozszerzoną wersję Twojego
\code{commit\_*}. Na początek wystarczą prostsze asercje SVA: „bańka nigdy nie ma
\code{mem\_write}”, „po stallu w ID jest ta sama instrukcja”.

\section{F. Oficjalne testy riscv-tests}
\emph{Około 1 wieczoru.} Testy \code{rv32ui-p-*} z repozytorium
\code{riscv-software-src/riscv-tests} używają tego samego protokołu TOHOST co my.
Wymagają toolchainu GNU i dostosowania skryptu linkera do naszej mapy pamięci.

\section{Lektura}

\begin{itemize}
  \item \textbf{Harris \& Harris}, \emph{Digital Design and Computer Architecture:
    RISC-V Edition}. Rozdziały 5, 7 i 8 to ścieżka tego kursu z innej perspektywy.
  \item \textbf{Patterson \& Hennessy}, \emph{Computer Organization and Design,
    RISC-V Edition}. Klasyka, rozdział 4 o potoku.
  \item \textbf{Hennessy \& Patterson}, \emph{Computer Architecture: A Quantitative
    Approach}. Następny poziom: wykonanie poza kolejnością, predykcja, pamięć.
  \item \textbf{The RISC-V Instruction Set Manual}, tomy I (Unprivileged) i II
    (Privileged). Zaskakująco czytelne, z uzasadnieniami decyzji projektowych.
  \item \textbf{Sutherland, Davidson, Flake}, \emph{SystemVerilog for Design}.
\end{itemize}

\section{Prawdziwe rdzenie do przeczytania}

\begin{center}
\small
\begin{tabularx}{\linewidth}{@{}l X@{}}
\toprule
\textbf{PicoRV32} & jeden plik, wielocyklowy, nastawiony na mały rozmiar \\
\textbf{SERV} & bit-serialny: 1 bit na takt, najmniejszy rdzeń RISC-V \\
\textbf{Ibex} (lowRISC) & 2--3 etapy, jakość produkcyjna, świetna dokumentacja weryfikacji \\
\textbf{VexRiscv} & konfigurowalny potok w SpinalHDL \\
\textbf{CVA6} & 64-bitowy, 6 etapów, uruchamia Linuksa \\
\bottomrule
\end{tabularx}
\end{center}

Porównaj ich decyzje ze swoimi: gdzie rozstrzygają skoki, jak zbudowany jest bank
rejestrów i jak obsługują pamięć.


\appendix
\chapter{Karta RV32I}
\label{app:karta}

Kodowanie wszystkich 37 instrukcji obliczeniowych, pamięciowych i sterujących RV32I
(bez \code{fence}, \code{ecall}, \code{ebreak}). Tabela jest wygenerowana z tablic
asemblera kursu (\plik{common/tools/rvtool.py}), który koduje instrukcje bit w bit
tak samo jak LLVM. Formaty i immediate'y: rys.~\ref{fig:formaty}.

% WYGENEROWANE z common/tools/rvtool.py — nie edytuj ręcznie
\begin{center}
\footnotesize
\renewcommand{\arraystretch}{1.0}
\begin{tabular}{@{}l c l l l l@{}}
\toprule
\textbf{Instr.} & \textbf{Fmt} & \textbf{opcode} & \textbf{funct3} & \textbf{funct7} & \textbf{Działanie} \\
\midrule
\code{lui} & U & \code{0110111} & \code{---} & \code{---} & rd = imm $\ll$ 12 \\
\code{auipc} & U & \code{0010111} & \code{---} & \code{---} & rd = pc + (imm $\ll$ 12) \\
\code{jal} & J & \code{1101111} & \code{---} & \code{---} & rd = pc+4; pc += imm \\
\code{jalr} & I & \code{1100111} & \code{000} & \code{---} & rd = pc+4; pc = (rs1+imm) \& \textasciitilde 1 \\
\midrule
\code{beq} & B & \code{1100011} & \code{000} & \code{---} & rs1 == rs2 \\
\code{bne} & B & \code{1100011} & \code{001} & \code{---} & rs1 != rs2 \\
\code{blt} & B & \code{1100011} & \code{100} & \code{---} & rs1 $<_s$ rs2 \\
\code{bge} & B & \code{1100011} & \code{101} & \code{---} & rs1 $\geq_s$ rs2 \\
\code{bltu} & B & \code{1100011} & \code{110} & \code{---} & rs1 $<_u$ rs2 \\
\code{bgeu} & B & \code{1100011} & \code{111} & \code{---} & rs1 $\geq_u$ rs2 \\
\midrule
\code{lb} & I & \code{0000011} & \code{000} & \code{---} & rd = sext(M[rs1+imm][7:0]) \\
\code{lh} & I & \code{0000011} & \code{001} & \code{---} & rd = sext(M[rs1+imm][15:0]) \\
\code{lw} & I & \code{0000011} & \code{010} & \code{---} & rd = M[rs1+imm] \\
\code{lbu} & I & \code{0000011} & \code{100} & \code{---} & rd = zext(M[rs1+imm][7:0]) \\
\code{lhu} & I & \code{0000011} & \code{101} & \code{---} & rd = zext(M[rs1+imm][15:0]) \\
\code{sb} & S & \code{0100011} & \code{000} & \code{---} & M[rs1+imm][7:0] = rs2[7:0] \\
\code{sh} & S & \code{0100011} & \code{001} & \code{---} & M[rs1+imm][15:0] = rs2[15:0] \\
\code{sw} & S & \code{0100011} & \code{010} & \code{---} & M[rs1+imm] = rs2 \\
\midrule
\code{addi} & I & \code{0010011} & \code{000} & \code{---} & rd = rs1 + imm \\
\code{slti} & I & \code{0010011} & \code{010} & \code{---} & rd = rs1 $<_s$ imm \\
\code{sltiu} & I & \code{0010011} & \code{011} & \code{---} & rd = rs1 $<_u$ imm \\
\code{xori} & I & \code{0010011} & \code{100} & \code{---} & rd = rs1 \^{} imm \\
\code{ori} & I & \code{0010011} & \code{110} & \code{---} & rd = rs1 | imm \\
\code{andi} & I & \code{0010011} & \code{111} & \code{---} & rd = rs1 \& imm \\
\code{slli} & I* & \code{0010011} & \code{001} & \code{0000000} & rd = rs1 $\ll$ shamt \\
\code{srli} & I* & \code{0010011} & \code{101} & \code{0000000} & rd = rs1 $\gg_u$ shamt \\
\code{srai} & I* & \code{0010011} & \code{101} & \code{0100000} & rd = rs1 $\gg_s$ shamt \\
\midrule
\code{add} & R & \code{0110011} & \code{000} & \code{0000000} & rd = rs1 + rs2 \\
\code{sub} & R & \code{0110011} & \code{000} & \code{0100000} & rd = rs1 - rs2 \\
\code{sll} & R & \code{0110011} & \code{001} & \code{0000000} & rd = rs1 $\ll$ rs2[4:0] \\
\code{slt} & R & \code{0110011} & \code{010} & \code{0000000} & rd = rs1 $<_s$ rs2 \\
\code{sltu} & R & \code{0110011} & \code{011} & \code{0000000} & rd = rs1 $<_u$ rs2 \\
\code{xor} & R & \code{0110011} & \code{100} & \code{0000000} & rd = rs1 \^{} rs2 \\
\code{srl} & R & \code{0110011} & \code{101} & \code{0000000} & rd = rs1 $\gg_u$ rs2[4:0] \\
\code{sra} & R & \code{0110011} & \code{101} & \code{0100000} & rd = rs1 $\gg_s$ rs2[4:0] \\
\code{or} & R & \code{0110011} & \code{110} & \code{0000000} & rd = rs1 | rs2 \\
\code{and} & R & \code{0110011} & \code{111} & \code{0000000} & rd = rs1 \& rs2 \\
\bottomrule
\end{tabular}
\end{center}


{\small I* --- format I, w którym \code{imm[11:5]} pełni rolę \code{funct7}, a
\code{imm[4:0]} to \code{shamt}. Wszystkie immediate'y poza U są rozszerzane
znakiem. \code{sext}/\code{zext}: rozszerzenie znakiem/zerami.}

\section*{Mapa pamięci maszyny kursu}

\begin{center}
\small
\begin{tabular}{@{}lll@{}}
\toprule
\code{0x0000\_0000}--\code{0x0000\_FFFF} & RAM 64 KiB & program, dane, stos \\
\code{0x1000\_0000} & TOHOST (W) & 1 = PASS, $(n \ll 1)|1$ = FAIL $n$ \\
\code{0x1000\_0004} & PUTCHAR (W) & znak \code{wdata[7:0]} \\
\code{0x1000\_0008} & CYCLES (R) & licznik cykli \\
\bottomrule
\end{tabular}
\end{center}

\chapter{Ściąga SystemVerilog}
\label{app:sv}

\section*{Składnia}
\begin{sv}
module nazwa import rv32i_pkg::*; #(parameter int W = 8) (
  input  logic         clk, rst,
  input  logic [W-1:0] a,
  output logic [W-1:0] y
);
  logic [W-1:0] t;                      // sygnał wewnętrzny
  localparam logic [3:0] STALA = 4'hA;  // stała lokalna

  assign t = a ^ {W{1'b1}};             // logika kombinacyjna (krótka)

  always_comb begin                     // logika kombinacyjna (dłuższa)
    y = '0;                             // wartość domyślna: brak zatrzasku
    case (t[1:0])
      2'd0:    y = a;
      default: y = t;
    endcase
  end

  always_ff @(posedge clk)              // rejestr
    if (rst) q <= '0;
    else     q <= y;
endmodule
\end{sv}

\section*{Literały i operatory}
\begin{center}
\small
\begin{tabularx}{\linewidth}{@{}l X@{}}
\toprule
\code{8'hFF}, \code{4'b1010}, \code{32'd7}, \code{'0}, \code{'1} & literały z rozmiarem; same zera/jedynki \\
\code{x[7:0]}, \code{x[8*i +: 8]} & wycinek stały; wycinek o zmiennym początku \\
\code{\{a, b\}}, \code{\{4\{b\}\}} & sklejanie; powielanie \\
\code{\& | \^{} \~{}} & bitowe; jako unarne: redukcja (\code{|x} = „jakikolwiek bit”) \\
\code{\&\& || !} & logiczne (1 bit) \\
\code{== !=}, \code{=== !==} & równość; dosłowna (z X i Z), tylko w testbenchach \\
\code{<< >> >>>} & przesunięcia; \code{>>>} arytmetyczne tylko dla signed \\
\code{\$signed(x)}, \code{\$unsigned(x)} & interpretacja ze znakiem / bez znaku \\
\code{W'(x)} & rzutowanie na szerokość \code{W} \\
\code{\$clog2(N)} & $\lceil \log_2 N \rceil$ (szerokość adresu dla $N$ elementów) \\
\bottomrule
\end{tabularx}
\end{center}

\section*{Zasady kursu}
\begin{enumerate}
  \item \code{always\_comb} + \code{=}, \code{always\_ff} + \code{<=}. Nigdy odwrotnie.
  \item W \code{always\_comb} każdy sygnał ma wartość domyślną na początku bloku.
  \item Każdy \code{case} ma \code{default}.
  \item Jeden zegar, reset synchroniczny aktywny w stanie 1.
  \item Stałe z pakietu, nie „magiczne liczby”.
  \item Uważaj na szerokości: \code{make lint} musi być czysty.
\end{enumerate}

\section*{Testbench}
\begin{sv}
`timescale 1ns/1ps
module tb_x;
  `include "tb_common.svh"
  logic clk = 0;
  always #5 clk = ~clk;
  initial begin
    tb_start();
    @(negedge clk);                       // zmieniaj wejścia na zboczu opadającym
    `CHECK_EQ(y, 8'h42, "opis")
    repeat (10) @(negedge clk);
    $display("a=%h b=%0d", a, b);
    tb_finish();
  end
endmodule
\end{sv}

\chapter{Ściąga poleceń}
\label{app:make}

\begin{center}
\small
\begin{tabularx}{\linewidth}{@{}l X@{}}
\toprule
\textbf{Polecenie} & \textbf{Działanie} \\
\midrule
\code{make test} & wszystkie testy modułu \\
\code{make unit [T=nazwa]} & testy jednostkowe (\plik{tb/tb\_nazwa.sv}) \\
\code{make wave T=nazwa} & test + przebiegi w GTKWave \\
\code{make lint} & Verilator \code{--lint-only -Wall} \\
\code{make synth T=moduł} & Yosys: komórki i najdłuższa ścieżka \\
\code{make progs [PAD=n] [SIM=verilator]} & 12 programów testowych na rdzeniu \\
\code{make prog P=08\_fib} & jeden program: wyjście, ślad \plik{build/P.trace}, listing \\
\code{make wave P=08\_fib} & przebiegi wykonania programu \\
\code{make cosim} & porównanie śladu z emulatorem \\
\code{make fuzz N=100 [LEN=300 SEED=1]} & losowe programy + porównanie \\
\code{make emu-test [P=...]} & (moduł 04) programy na emulatorze \\
\code{make run P=hello [CORE=08]} & (moduł 07) program w C na rdzeniu \\
\code{make clean} & usuwa \plik{build/} \\
\bottomrule
\end{tabularx}
\end{center}

\section*{Narzędzia w \plik{common/tools}}
\begin{sh}
python3 rvtool.py asm prog.S -o prog.hex -l prog.lst [-I katalog] [--pad-nops N]
python3 rvtool.py dis prog.hex          # albo: dis 0x00a50533
python3 rvtool.py bin2hex prog.bin -o prog.hex
python3 tracecmp.py ref.trace dut.trace
python3 rvgen.py --seed 7 -n 300 -o rand.S [--mext]
\end{sh}

\chapter{Słowniczek}
\label{app:slownik}

\begin{center}
\small
\begin{tabularx}{\linewidth}{@{}l l X@{}}
\toprule
\textbf{Polski} & \textbf{Angielski} & \textbf{Znaczenie} \\
\midrule
przerzutnik & flip-flop (FF) & element pamiętający bit na zboczu zegara \\
zatrzask & latch & element pamiętający sterowany poziomem; w RTL zwykle błąd \\
ścieżka krytyczna & critical path & najdłuższa droga kombinacyjna między rejestrami \\
bank rejestrów & register file & 32 rejestry z portami odczytu i zapisu \\
ścieżka danych & datapath & bloki przetwarzające dane (ALU, rejestry, multipleksery) \\
jednostka sterująca & control unit & dekoder wytwarzający sygnały sterujące \\
potok & pipeline & podział wykonania na etapy oddzielone rejestrami \\
hazard & hazard & sytuacja, w której potok dałby zły wynik bez dodatkowej logiki \\
przekazywanie & forwarding / bypass & dostarczenie wyniku prosto z rejestru potoku \\
wstrzymanie & stall & zatrzymanie części potoku na takt \\
bańka & bubble & pusty etap potoku (\code{valid = 0}) \\
unieważnienie & flush & zamiana instrukcji z błędnej ścieżki w bańki \\
zatwierdzenie & commit / retire & zakończenie instrukcji, widoczne w śladzie \\
stała natychmiastowa & immediate & stała zakodowana w instrukcji \\
rozszerzenie znakiem & sign extension & kopiowanie bitu znaku na starsze bity \\
co-symulacja & co-simulation & porównywanie RTL z modelem w trakcie wykonania \\
model referencyjny & golden model & niezależna implementacja traktowana jako wzorzec \\
pokrycie & coverage & miara tego, które zachowania zostały sprawdzone \\
\bottomrule
\end{tabularx}
\end{center}


\end{document}
